\documentclass[10pt, a4paper]{iopart}
% \usepackage{iopams} 

% REVISION:
% DELETE ALL LINES CONTAINING
%	\REV
% TO REMOVE THE HIGHLIGHTING
% OR COMMENT NEXT LINE
% \usepackage{color}\newcommand{\REV}[1]{{\color{magenta}#1}}

% COMMENT THIS LINE FOR PRODUCTION:
% \usepackage[a5paper,margin=1cm]{geometry}
% \usepackage{a4wide}

% IOP STYLEFILE COMPATIBILITY HACK 
\expandafter\let\csname equation*\endcsname\relax
\expandafter\let\csname endequation*\endcsname\relax

% 
\usepackage{amsmath, amsthm}
% \usepackage{amssymb}
\allowdisplaybreaks

% Load readable font (avoiding conflict with amssymb)
\makeatletter
\@ifpackageloaded{amssymb}{}{\usepackage[bitstream-charter]{mathdesign}}
\makeatother

\usepackage{color}

\usepackage[pdftex, bookmarksopen, pagebackref,
            pdftitle={Tikhonov and Landweber convergence rates: characterization by interpolation spaces},
            pdfauthor={R. Andreev},
            colorlinks, linkcolor=black, urlcolor=black, citecolor=black,
]{hyperref}


\newcommand{\IR}{\mathbb{R}}
\newcommand{\IN}{\mathbb{N}}
\newcommand{\IT}{\mathbb{T}}
\newcommand{\IX}{\mathbb{X}}
\newcommand{\argmax}{\mathop{\mathrm{argmax}}}
\newcommand{\norm}[2]{     \| #1       \|_{ #2 }}
\newcommand{\Norm}[2]{\left\| #1 \right\|_{ #2 }}
\newcommand{\normiii}[2]{\vert\kern-0.25ex\vert\kern-0.25ex\vert #1 \vert\kern-0.25ex \vert\kern-0.25ex\vert_{ #2 }}
\newcommand{\Normiii}[2]{\left\vert\kern-0.25ex\left\vert\kern-0.25ex\left\vert #1 \right\vert\kern-0.25ex\right\vert\kern-0.25ex\right\vert_{ #2 }}
\newcommand{\scalar}[2]{     \langle #1       \rangle_{ #2 }}
\newcommand{\Scalar}[2]{\left\langle #1 \right\rangle_{ #2 }}

\newtheorem{lemma}{Lemma}[section]
\newtheorem{proposition}[lemma]{Proposition}
\newtheorem{theorem}[lemma]{Theorem}
\newtheorem{corollary}[lemma]{Corollary}
\newtheorem{remark}[lemma]{Remark}
\theoremstyle{remark}
\theoremstyle{definition}
\newtheorem{assumption}[lemma]{Assumption}
\newtheorem{definition}[lemma]{Definition}
\newtheorem{example}[lemma]{Example}

\newcommand{\TODO}[1]{\textrm{\color{red}TODO: #1}}
\definecolor{someblue}{rgb}{0,0,.8}
\newcommand{\ra}[1]{{\iftoggle{preview}{}{\color{someblue}}#1}}

\providecommand{\REV}[1]{{#1}}
\newcommand\sectiom{\vfil\penalty-9999\vfilneg\section}
\usepackage{suffix}
\WithSuffix\newcommand\sectiom*{\vfil\penalty-9999\vfilneg\section*}

% NOTES:
%
% Gilyazov/Goldman: alpha-processes
% Scherzer/Neubauer: alpha-processes for nonlinear
% Barzilai/Borwein: http://imajna.oxfordjournals.org/content/8/1/141.abstract


%%%%%%%%%%%%%%%%%%%%%%%%%%%%%%%%%%%%%%%%%%%%%%%%%%%%%%%%%%%%%%%%%%%%%%%%%%%%%%%%
\begin{document}
%%%%%%%%%%%%%%%%%%%%%%%%%%%%%%%%%%%%%%%%%%%%%%%%%%%%%%%%%%%%%%%%%%%%%%%%%%%%%%%%

\title%
[\hfill Tikhonov and Landweber rates by interpolation spaces \hfill]%
{Tikhonov and Landweber convergence rates: characterization by interpolation spaces}

\author{R.~Andreev}

\address{%
	Univ Paris Diderot, Sorbonne Paris Cit\'e, LJLL (UMR 7598 CNRS), F-75205 Paris, France
}
\ead{roman.andreev@upmc.fr}
\vspace{10pt}
\begin{indented}
\item[]\today
\end{indented}


\begin{abstract}
	Algebraic convergences rates 
	of (iterated) Tikhonov regularization
	for linear inverse problems in Hilbert spaces
	are 
	characterized by 
	the membership of 
	the exact solution 
	to
	intermediate spaces
	produced by the K-method of real interpolation.
	Similar results are obtained for the Landweber iteration.
\end{abstract}

% \pacs{00.00, 20.00, 42.10}

\vspace{2pc}
\noindent{\it Keywords}: 
Tikhonov, Landweber, regularization, convergence rates, sharp converse results,
interpolation spaces, K-functional,
47A52, % General theory of linear operators / Ill-posed problems, regularization
65F10, % Numerical linear algebra / Iterative methods for linear systems
65R30, % Integral equations, integral transforms / Improperly posed problems
65J20, % Numerical analysis in abstract spaces / Improperly posed problems; regularization
65J22  % Numerical analysis in abstract spaces / Inverse problems

% \maketitle


%%%%%%%%%%%%%%%%%%%%%%%%%%%%%%%%%%%%%%%%%%%%%%%%%%%%%%%%%%%%%%%%%%%%%%%%%%%%%%%%
\section{Introduction}
\label{s:0}
%%%%%%%%%%%%%%%%%%%%%%%%%%%%%%%%%%%%%%%%%%%%%%%%%%%%%%%%%%%%%%%%%%%%%%%%%%%%%%%%

It has long been known
that 
the H\"older type 
source conditions 
% \cite[(3.35)]{EnglHankeNeubauer1996}
originally used
to obtain algebraic convergence rates
for Tikhonov regularization
are only sufficient,
not necessary.
%
Consequently,
other concepts 
have been introduced 
to completely characterize those rates
%
%
such as 
the spectral decay condition of \cite[Theorem 2.1]{Neubauer1997},
distance functions \cite[p.3]{FlemmingHofmannMathe2011}, 
and
variational source conditions 
\cite{Flemming2012};
% Flemming{2013}
pointers to the origins of those concepts can be found therein.
%
%
We propose here a more concise condition
through 
the notion of interpolation spaces
and 
establish links 
to the concepts just listed.
%
Specifically,
we argue that
the intermediate spaces
$(E_0, E_1)_{\theta, q}$
produced by the $K$-method of real interpolation
\cite{BerghLofstrom1976, BennettSharpley1988}
with fine index $q = \infty$
naturally capture 
the essential behavior
of (iterated) Tikhonov regularization,
that is:
convergence rates, converse results, and saturation,
in the noise-free and the noisy case.
%
This is not unexpected, given 
the resemblance of the $K$-functional \eqref{e:K}
and 
the Tikhonov functional \eqref{e:J},
but we systematically quantify this connection
by careful estimates
with particular attention to the limiting cases
$\theta = 0, 1$.
%
In a similar vein, 
the relationship between near-minimizers for ``$L$-functionals''
and
Tikhonov regularization
was highlighted
in \cite[Chapter 6]{KislyakovKruglyak2013}.
%
%
%
We prove analogous convergence and converse results for 
the Landweber iteration,
%in terms of membership of the exact solution
%to such intermediate spaces,
and comment on the applicability of the discrepancy principle
as a stopping rule.


%

This note consists of two main parts.
%
In the first part (Section \ref{s:i})
we develop the required preliminaries
of the $K$-method of real interpolation,
introduce the different intermediate subspaces,
and describe them and their interrelations using 
spectral theory in Hilbert spaces.
%
In passing, we relate
to the concepts of distance functions 
and 
variational source conditions.
%
%
In the second part (Section \ref{s:ip})
we elaborate on how
the convergence rates of
Tikhonov regularization
and
Landweber iteration
are characterized 
in terms of the intermediate subspaces.

%

We write $A \lesssim B$
to mean that there is a constant $C \geq 0$ 
independent of the parameters specified by quantifiers
such that $A \leq C B$.
%
If, in addition, $B \lesssim A$, we write $A \sim B$.


%%%%%%%%%%%%%%%%%%%%%%%%%%%%%%%%%%%%%%%%%%%%%%%%%%%%%%%%%%%%%%%%%%%%%%%%%%%%%%%%
\section{Preliminaries}
\label{s:i}
%%%%%%%%%%%%%%%%%%%%%%%%%%%%%%%%%%%%%%%%%%%%%%%%%%%%%%%%%%%%%%%%%%%%%%%%%%%%%%%%

%%%%%%%%%%%%%%%%%%%%%%%%%%%%%%%%%%%%%%%%%%%%%%%%%%%%%%%%%%%%%%%%%%%%%%%%%%%%%%%%
\subsection{Interpolation spaces}
%%%%%%%%%%%%%%%%%%%%%%%%%%%%%%%%%%%%%%%%%%%%%%%%%%%%%%%%%%%%%%%%%%%%%%%%%%%%%%%%

Let $X$ be a Banach space (here and henceforth: over the reals).
%
Let $X_1 \subset X$ be another Banach space,
continuously 
%\TODO{and densely?} 
embedded in $X$.
%
We write $\norm{\cdot}{0} := \norm{\cdot}{}$ and $\norm{\cdot}{1}$
for the norms of $X$ and $X_1$, respectively.
%
%
The $K$-functional
is defined as
\begin{align}
\label{e:K}
	K_t(x) :=
	\inf_{x_1 \in X_1}
	(
		\norm{x - x_1}{0}^2
		+
		t^2
		\norm{x_1}{1}^2
	)^{1/2}
	,
	\quad x \in X,
	\quad t > 0.
\end{align}
%
For real $0 < \theta < 1$ and $1 \leq q \leq \infty$
the $K$-method of real interpolation
defines intermediate subspaces 
$X_1 \subset (X, X_1)_{\theta, q} \subset X$
based on the integrability of $t \mapsto K_t(x)$.
%
Here we are only interested in the case $q = \infty$ 
with one of the spaces embedded into the other,
and therefore
refer to standard sources
such as 
\cite{BerghLofstrom1976, BennettSharpley1988}
for general definitions.
%
For any $0 \leq \theta \leq 1$ and any $x \in X$ set
\begin{align}
\label{e:||nu}
	\norm{x}{\theta:1}
	:=
	\sup_{t > 0}
	t^{-\theta}
	K_t(x)
	,
\end{align}
possibly infinite.
%
We point out
that 
we include the limiting cases $\theta = 0$ and $\theta = 1$
in this definition.
%
The reason for the unusual notation 
will become apparent in Section \ref{s:i:H}
where 
instead of $X_1$
we will 
consider a family of subspaces $X_\gamma \subset X$
parametrized by $\gamma$.
%
Define the spaces
%the set
\begin{align}
	X_{\theta:1} 
	:=
	( X, X_1 )_{\theta, \infty}
	:=
	\{
		x \in X :
		\norm{x}{\theta:1} < \infty
	\}
\end{align}
with the norm $\norm{\cdot}{\theta:1}$.
%
For $0 < \theta < 1$,
these are Banach spaces.
%
Moreover,
the following embeddings 
$
	X_1 \subset X_{\theta:1} \subset X
$
are continuous.
%
%
The space $X_{\theta:1}$ need not coincide with $X_\theta$
for $\theta = 0,1$,
see the remarks on the Gagliardo completion in \cite[Chapter 5, \S1]{BennettSharpley1988},
but it will be the case in 
the more specific setting of Section \ref{s:i:H}.



%%%%%%%%%%%%%%%%%%%%%%%%%%%%%%%%%%%%%%%%%%%%%%%%%%%%%%%%%%%%%%%%%%%%%%%%%%%%%%%%
\subsection{The constant \texorpdfstring{$N_\theta$}{Nnu}}
%%%%%%%%%%%%%%%%%%%%%%%%%%%%%%%%%%%%%%%%%%%%%%%%%%%%%%%%%%%%%%%%%%%%%%%%%%%%%%%%

For $0 \leq \theta \leq 1$, 
the constant $1 \leq N_\theta \leq \sqrt{2}$ given by
\begin{align}
\label{e:N=sup}
	N_{\theta}^{-2}
	:=
	\theta^\theta (1-\theta)^{1-\theta}
	=
	\sup_{\lambda \in [0, 1]}
	\lambda^\theta (1 - \lambda)^{1 - \theta}
	=
	\sup_{s > 0}
	\frac{s^{2(1 - \theta)}}{s^2 + 1}
\end{align}
will play a recurrent role.
%
By convention, $N_0 := N_1 := 1$,
making $\theta \mapsto N_\theta$ continuous on $[0, 1]$.
%
For example,
if
$a, b > 0$
then
Young's inequality with exponents $p = 1/(1-\theta)$ and $q = 1 / \theta$
gives
\begin{align}
\label{e:Nab}
	N_{\theta}^2
	a^{1-\theta} b^{\theta} 
	=
	(a / (1-\theta))^{1-\theta} (b / \theta)^{\theta}
	\leq
	a + b
	.
\end{align}
%
As another example,
for any real $\lambda \geq 0$ and $t > 0$,
any real $k \geq 1/2$, and any $0 \leq \theta \leq 1$,
\begin{align}
\label{e:Ncomm}
	N_\theta^{2 (1 - 2 k)} 
	\frac{t^{2 (1 - \theta)}}{t^2 + \lambda^{2 k}}
	\leq
	\left(
		\frac{\alpha^{1 - \theta}}{\alpha + \lambda}
	\right)^{2k}
	%
	\quad\text{for}\quad
	%
	\alpha 
	=
	t^{1/k} \left( \frac{1 - \theta}{\theta} \right)^{1 - 1/(2k)}
	.
\end{align}
%
Indeed, after algebraic simplification, 
the inequality in
\eqref{e:Ncomm}
is equivalent to
$
	(\alpha + \lambda)^p
	\leq
	(1-\theta)^{1-p} \alpha^p + \theta^{1-p} \lambda^p
$,
itself a consequence of H\"older's inequality
with $p = 2k \geq 1$ as one of the exponents.

%In view of \eqref{e:ii} and \eqref{e:|Snu|},
%it is tempting to normalize by $N_\theta$ in the definition \eqref{e:||nu}
%of the interpolation norm
%but
%we chose not to do so
%to avoid the proliferation of the constant in certain key statements
%such as \eqref{e:ga-nu} or \eqref{e:p:tikk}.


\subsection{Distance function}
\label{s:i:dr}

Distance functions were introduced
as a means
to characterize the regularization error of linear regularization operators
in \cite{FlemmingHofmannMathe2011},
previously also in \cite[Theorem 2.12]{BakushinskyKokurin2004}.
%
We briefly comment on the relation to the $K$-functional.
%
Let $0 < \theta < 1$.
%
Fix $x \in X$.
%
Define
the distance function
\begin{align}
\label{e:dr} 
	d(r) 
	:=
	\inf
	\{ \norm{ x - x_1 }{0} : x_1 \in X_1, \; \norm{x_1}{1} \leq r \}
	,
	\quad
	r \geq 0
	.
\end{align}
%
This function is nonnegative,
nonincreasing,
bounded by $d(0) = \norm{x}{0}$,
and convex.
%
% It achieves $d(r) = 0$ at a finite $r \geq 0$
% if and only if $x \in X_1$. \TODO{Really?}
%
% Moreover, $d(\infty) = 0$ if and only if $X_1$ is dense in $X$.
%
% Therefore, once it is constant, it stays constant.
%
To simplify the notation,
we shall write $\norm{x_1}{1} \leq r$, or similar,
without mentioning that $x_1 \in X_1$.
%
It is clear the $d$ has compact support
if and only if $x \in X_1$.

Inspection of the definitions 
of the $K$-functional \eqref{e:K} and the distance function \eqref{e:dr}
reveals
that
\begin{align}
\label{e:Kd*}
	K_t(x) = \inf_{r > 0} (d(r)^2 + t^2 r^2)^{1/2}
	\sim
	\inf_{r > 0} (d(r) + t r)
%	=
%	-\sup_{r > 0} ((-t) r  - d(r))
	=
	-d^*(-t)
	,
\end{align}
where $d^*$ is the Legendre--Fenchel conjugate of $d$.
%
Thus the distance function \eqref{e:dr} 
characterizes the subspace $X_{\theta:1} \subset X$.
%
The first characterization,
boundedness
of $|t^{-\theta} d^*(-t)|$,
is obvious from \eqref{e:||nu} and \eqref{e:Kd*}.
%
The second characterization
is the behavior of $d$ at infinity,
more precisely the identity
\begin{align}
\label{e:D=E}
	E %_\theta(x)
	=
	N_{\theta}^{-1}
	\norm{x}{\theta:1}
	=
	D %_\theta(x)
	,
\end{align}
%whenever either is finite,
where
\begin{align}
\label{e:ED}
 	E%_\theta(x)
	:=
	\sup_{r > 0}
	\inf_{ \norm{x_1}{1} \leq r }
	\norm{ x - x_1 }{0}^{1-\theta} 
	\norm{x_1}{1}^\theta
	\quad\text{and}\quad
 	D%_\theta(x) 
	:=
	\sup_{r > 0}
	d(r)^{1 - \theta} 
	r^\theta
	.
\end{align}
%
%
%
From \eqref{e:D=E} one infers
the more qualitative observation
that
$x \in X_{\theta:1}$
if and only if
the distance function \eqref{e:dr} exhibits
the asymptotic decay rate
$d(r) = \mathcal{O}(r^{-\theta / (1-\theta)})$ for $r \to \infty$.
%

\begin{proof}[{Proof of \eqref{e:D=E}}]
	The proof 
	% of \eqref{e:D=E} 
	is in the three steps:
	%
	
	a) $E \leq N_{\theta}^{-1} \norm{x}{\theta:1}$.
	%
% 	Set $p := 1 / (1 - \theta)$ and $q := 1 / \theta$.
	%
	Estimating the $\inf$ of $E$ as in \eqref{e:Nab}
	for $a := \norm{ x - x_1 }{0}^2$ and $b := t^2 \norm{x_1}{1}^2$,
	\begin{align*}
		\inf_{\norm{x_1}{1} \leq r} (\cdots)
	\leq
		N_{\theta}^{-1}
		t^{-\theta}
		\inf_{\norm{x_1}{1} \leq r}
		(
			\norm{x - x_1}{0}^2
			+
			t^2 \norm{x_1}{1}^2
		)^{1/2}
	\leq
		N_{\theta}^{-1}
		\sup_{t > 0}
		t^{-\theta}
		K_t(x)
% 	=
% 		N_{\theta}^{-2}
% 		\norm{x}{\theta}
		.
	\end{align*}
	%
	Here,
	the condition $\norm{x_1}{1} \leq r$ is redundant
	if $t$ is large enough,
	so the infimum becomes $K_t(x)$,
	explaining the second inequality.
	%
	Taking the supremum over $r > 0$ proves the claim.
	
	b) $\norm{x}{\theta:1} \leq N_{\theta} D$.
	%
	The definition \eqref{e:ED} of $D$ 
	implies
	$d(r) \leq (r^\theta / D)^{1/(\theta - 1)}$.
	%
	Using this in \eqref{e:Kd*},
	then
	computing the infimum
	yields
	$K_t(x) \leq t^\theta N_{\theta} D$.
	%
	Now
	multiply by $t^{-\theta}$
	and take $\sup_{t > 0}$.
	
	
	c) $D \leq E$.
	%
	If $D$ is finite then $d(\infty) = 0$.
	%
	Given that $d$ is convex,
	$d$ is strictly decreasing
	(unless where it vanishes).
	%
	Thus,
	for each $r > 0$,
	the infimum in $d(r)$
	is achieved at $\norm{x_1}{1} = r$.
	%
	Concerning $E$,
	we may suppose 
%	without loss of generality 
	that 
	the $\sup \inf$ is assumed at $\norm{x_1}{1} = r$,
	adjusting $r$ if necessary.
	%
	Hence, both $D$ and $E$ equal
%	the same quantity
%	which is
	$
		\sup_{r > 0} 
		\inf_{ \norm{x_1}{1} = r }
		\norm{ x - x_1 }{0}^{1-\theta} 
		\norm{x_1}{1}^\theta
	$.
	%
	%
	This establishes \eqref{e:D=E}.
\end{proof}
%

%



\subsection{Interpolation inequality of operators}

Let $Y$ be a Banach space.
%
Let $S : X \to Y$ be a bounded linear operator with norm $C_0 \geq 0$.
%
Assume that 
$S|_{X_1} : X_1 \to Y$ 
is also a bounded linear operator, with norm $C_1 \geq 0$.
%
Then, for any $0 < \theta < 1$, 
\begin{align}
\label{e:|Snu|}
% 	(\text{\color{red} $N_\theta$ here?})
	\norm{S x}{Y}
	\leq
	N_{\theta}
	C_0^{1-\theta} C_1^\theta
	\norm{x}{\theta:1}
	\quad
	\forall x \in X_{\theta:1}
	.
\end{align}
%
%Thus,
%$S|_{X_{\theta:1}} : X_{\theta:1} \to Y$
%is also bounded.
%
Indeed,
% in
% the nontrivial case
% $0 < \theta < 1$
let $x \in X_{\theta:1}$.
%
Given any $x_1 \in X_1$, 
write $x = (x - x_1) + x_1$,
apply the triangle inequality with boundedness of $S$,
and estimate by Cauchy--Schwarz:
$ %\begin{align}
	\norm{S x}{Y}
% 	\leq
% 	\norm{S (x - x_1)}{Y}
% 	+
% 	\norm{S x_1}{Y}
	\leq
	C_0 \norm{x - x_1}{0}
	+
	C_1 \norm{x_1}{1}
	\leq
	t^\theta (C_0^2 + t^{-2} C_1^2)^{1/2}
	\times
	t^{-\theta}
	(\norm{x - x_1}{0}^2 + t^2 \norm{x_1}{1}^2)^{1/2}
	.
$ %\end{align}
%
Inserting 
an
infimum over $x_1 \in X_1$ 
then a supremum over $t > 0$
in the second factor,
in view of \eqref{e:||nu}
we obtain
$ %\begin{align}
	\norm{S x}{Y}
	\leq
	t^\theta (C_0^2 + t^{-2} C_1^2)^{1/2}
	\times
	\norm{x}{\theta:1}
	.
$ %\end{align}
%
Minimization over $t > 0$ gives \eqref{e:|Snu|}.





%\subsection{Interpolation inequality of norms}
%
%For any $0 \leq \theta \leq 1$,
%the interpolation inequality
%\begin{align}
%\label{e:ii}
%	\norm{x}{\theta:1}
%	\leq
%	N_\theta^{-1}
%	\norm{x}{0}^{1 - \theta}
%	\norm{x}{1}^\theta
%	\quad
%	\forall x \in X_1
%\end{align}
%is valid.
%%
%To prove this for the nontrivial case $0 < \theta < 1$,
%we use $\norm{x}{\theta:1} = N_{\theta} E$ from \eqref{e:ED},
%estimate above
%the $\sup \inf$ in $E$
%by
%$\sup_{\lambda \in [0, 1]}$ 
%while choosing $x_1 := \lambda x$,
%and apply \eqref{e:N=sup}.



%%%%%%%%%%%%%%%%%%%%%%%%%%%%%%%%%%%%%%%%%%%%%%%%%%%%%%%%%%%%%%%%%
\subsection{A lemma for measures}
\label{s:}
%%%%%%%%%%%%%%%%%%%%%%%%%%%%%%%%%%%%%%%%%%%%%%%%%%%%%%%%%%%%%%%%%

We will 
call a nonnegative finite measure $\mu$ on Borel subsets of $[0, \infty)$
a ``Borel measure on $[0, \infty)$''.
%
For any such $\mu$ and any real $\nu \geq 0$
we define
$\normiii{\mu}{\nu}$
by
\begin{align}
	\normiii{\mu}{\nu}^2
	:=
	\sup_{t > 0} t^{-2\nu} \mu([0, t))
	.
\end{align}


The following Lemma 
records
two useful properties of $\normiii{\cdot}{\nu}$.

\begin{lemma}
\label{l:mu}
	Let $\mu$ be a Borel measure on $[0, \infty)$.
	%
	Let $\nu > 0$.
	%
	Then
	\begin{align}
	\label{e:Tail2}
		\mu([0, \Lambda))
		+
		\Lambda^{2 \gamma}
		\int_{[\Lambda, \infty)}
		\lambda^{-2 \gamma}
		d\mu(\lambda)
		\leq
		\frac{\gamma}{\gamma - \nu}
		\Lambda^{2 \nu}
		\normiii{\mu}{\nu}^2
		\quad
		\forall \Lambda > 0
		\quad
		\forall \gamma > 0
		,
	\end{align}
	and
	\begin{align}
	\label{e:Tail1}
		\int_{[0, \Lambda)} \lambda^{-2 r} d\mu(\lambda)
	\leq
		\frac{r + \nu}{\nu}
		\Lambda^{2 \nu}
		\normiii{\mu}{r + \nu}^2
	\quad
		\forall \Lambda > 0
		\quad
		\forall r \geq 0
		,
	\end{align}
	whenever the right-hand-side is finite.
\end{lemma}

\begin{proof}
	Define the left-continuous function
	$I(\Lambda) := \mu([0, \Lambda))$
	for $\Lambda \geq 0$.
	%
	Fix $\Lambda > 0$.
	%
	Writing the integral as a Riemann--Stieltjes integral, 
	and 
	integrating by parts we have
	\begin{align}
		\int_{[\Lambda, \infty)}
		\lambda^{-2 \gamma}
		d\mu(\lambda)
		=
		\int_{[\Lambda, \infty)}
		\lambda^{-2 \gamma}
		d I(\lambda)
		=
		-\Lambda^{-2 \gamma} I(\Lambda) 
		+
		2 \gamma
		\int_{[\Lambda, \infty)}
		\lambda^{-2 \gamma - 1}
		I(\lambda)
		d \lambda
		.
	\end{align}
	%
	Estimating 
	$I(\lambda) \leq \lambda^{2 \nu} \normiii{\mu}{\nu}^2$ 
	under the integral,
	evaluating,
	and rearranging
	leads to \eqref{e:Tail2}.
	%
	Similarly,
	\begin{align}
		\int_{[0, \Lambda)}
		\lambda^{-2 r}
		d \mu(\lambda)
	\leq
		\int_{[0, \Lambda]} 
		\lambda^{-2 r}
		d I(\lambda)
	=
		\left. 
			\lambda^{-2 r} I(\lambda) 
		\right|_{\lambda = 0}^{\lambda \searrow \Lambda}
		+
		2 r
		\int_{[0, \Lambda]}
		\lambda^{-2 r - 1}
		I(\lambda)
		d \lambda
		.
	\end{align}
	%
	Estimating $I(\lambda) \leq \lambda^{2 (r + \nu)} \normiii{\mu}{r + \nu}^2$
	and evaluating the integral yields \eqref{e:Tail1}.
\end{proof}

{%
% \footnotesize
\begin{proof}[Alternative proof] % of \eqref{e:Tail2}--\eqref{e:Tail1}]
	Fix $\Lambda > 0$.
	%
	%
	%
	Ad \eqref{e:Tail2}:
	%
	Consider
	$A_s := [\Lambda, \infty) \cap \{ \lambda \geq 0 : \lambda^{-2 \gamma} > s \}$.
	%
	If $s \geq \Lambda^{-2 \gamma}$ then $A_s$ is empty,
	otherwise
	$A_s = [\Lambda, s^{-1/(2\gamma)})$.
	%
	Therefore
	\begin{align}
		\int_{[\Lambda, \infty)} \lambda^{-2 \gamma} d\mu(\lambda)
		=
		\int_0^{\Lambda^{-2 \gamma}} \mu(A_s) ds
		=
		-\Lambda^{-2\gamma} \mu([0, \Lambda))
		+
		\int_0^{\Lambda^{-2 \gamma}} \mu([0, s^{-1/(2\gamma)})) ds
	\end{align}
	%
	Estimating $\mu([0, t)) \leq t^{2 \nu} \normiii{\mu}{\nu}^2$ under the integral
	and evaluating
	yields \eqref{e:Tail2}.
	%
	%
	%
	Ad \eqref{e:Tail1}:
	%
	The statement is trivial for $r = 0$, so suppose $r > 0$.
	%
	Consider 
	$B_s := [0, \Lambda) \cap \{ \lambda \geq 0 : \lambda^{-2 r} > s \}$.
	%
	If $s \leq \Lambda^{-2r}$ 
	then $B_s = [0, \Lambda)$ 
	and 
	$\mu(B_s) \leq \Lambda^{2 (\nu + r)} \normiii{\mu}{\nu + r}^2$.
	%
	Otherwise
	$B_s = [0, s^{-1/(2r)})$ ,
	so that
	$\mu(B_s) \leq s^{-(\nu + r)/r} \normiii{\mu}{\nu + r}^2$.
	%
	Using this in 
	$
		\int_{[0, \Lambda)} \lambda^{-2 r} d\mu(\lambda)
	=
		\int_0^\infty \mu(B_s) ds
	$
	yields \eqref{e:Tail1}.
	%
% 	Therefore
% 	\begin{align}
% % 	\textstyle
% 		\int_{[0, \Lambda)} \lambda^{-2 r} d\mu(\lambda)
% 	=
% 		\int_0^\infty \mu(B_s) ds
% 	\leq
% 		\left(
% 			\int_0^{\Lambda^{-2 r}} \Lambda^{2 (\nu + r)} ds
% 			+
% 			\int_{\Lambda^{-2 r}}^\infty s^{-(\nu + r) / r} ds
% 		\right)
% 		\normiii{\mu}{\nu + r}^2
% 		,
% 	\end{align}
% 	which evaluates to \eqref{e:Tail1}.
\end{proof}
}
	
%%%%%%%%%%%%%%%%%%%%%%%%%%%%%%%%%%%%%%%%%%%%%%%%%%%%%%%%%%%%%%%%%
\subsection{Spectral theory in Hilbert spaces}
\label{s:i:H}
%%%%%%%%%%%%%%%%%%%%%%%%%%%%%%%%%%%%%%%%%%%%%%%%%%%%%%%%%%%%%%%%%

Suppose that $X$ and $Y$ are real Hilbert spaces.
%
Let $T : X \to Y$ be a nonzero bounded linear operator.
%
Replacing $X$ by $X / \ker T$ if necessary
(with the usual quotient norm),
we assume that
\begin{align}
\label{e:inj}
	\text{$T$ is injective}
	.
\end{align}
%
Let $E$ denote the projection valued spectral measure of $T^* T$.
%
Then $E$ is compactly supported in $[0, \infty)$ since $T$ is bounded,
and injectivity of $T$ is equivalent to $E_{\{0\}} = 0$,
since 
$\mathop{\mathrm{range}}(E_{\{\lambda\}}) = \ker( T^* T - \lambda I)$.
%
For $x \in X$ we define 
the Borel measure $\mu_x$ on $[0, \infty)$
by
$\mu_x(A) := \norm{E_A x}{}^2$
and
the associated quantity
\begin{align}
\label{e:normiii}
	\normiii{x}{\nu} 
	:= 
	\normiii{\mu_x}{\nu}
	=
	\sup_{t > 0}
	t^{-\nu}
	\norm{ E_{[0, t)} x }{}
	,
	\quad
	\nu \geq 0
	.
\end{align}
%
The subset 
\begin{align}
	\IX_\nu :=
	\{
		x \in X:
		\normiii{x}{\nu} < \infty
	\}
\end{align}
of $X$
%\TODO{mod $\ker T^* T$?}
is indeed a Banach space 
equipped with the norm
$\normiii{\cdot}{\nu}$.
%
A description of this subspace 
as an interpolation space
is subsequently given 
in 
Proposition \ref{p:ga-nu}.
%
%
The definition of $\IX_\nu$ is inspired by 
the work \cite{Neubauer1997}.


%

For all real $\gamma \geq 0$,
we define the Banach space $X_\gamma \subset X$ 
as
\begin{align}
\label{e:Xga}
	X_\gamma := \mathop{\mathrm{range}} ( (T^* T)^\gamma )
	%
	\quad\text{with the norm}\quad
	%
	\norm{x}{\gamma}
	:=
	\norm{ (T^* T)^{-\gamma} x }{},
%	:=
%	\inf
%	\{
%		\norm{\omega}{}
%		:
%		\omega \in X, \;
%		(T^* T)^\gamma \omega = x
%	\},
	\quad
	x \in X_\gamma
	.
\end{align}
%
%
% By injectivity \eqref{e:inj} of $T$,
% there is only one candidate for the infimum.
%
%
%The spaces $X_\gamma$, $\gamma \geq 0$, in fact form part of a Hilbert scale.
%
In terms of the spectral measure $E$, we can write
(note $E_{\{0\}} = 0$)
\begin{align}
\label{e:0ga}
	\norm{x}{0}^2
	=
	\int_{[0, \infty)} d \mu_{x}(\lambda)
	\quad\text{and}\quad
	\norm{x}{\gamma}^2
	=
	\int_{[0, \infty)} \lambda^{-2\gamma} d \mu_{x}(\lambda)
	\quad
	\forall
	x \in X_\gamma
	.
\end{align}
%


%

%

For any real $0 \leq \nu \leq \gamma$,
we define
the interpolation space
\begin{align}
	X_{\nu : \gamma}
	:=
	( X, X_\gamma )_{\nu/\gamma, \infty}
	\quad\text{with norm}\quad
	\norm{\cdot}{\nu : \gamma}
	,
\end{align}
and will denote the corresponding $K$-functional by $K_t^\gamma$.
%
One can check
(most easily in the case that $T$ is compact)
that
for any $x \in X$ and $t > 0$,
\begin{align}
\label{e:KtI}
	| K_t^\gamma(x) |^2
=
%	\inf_{x_1 \in X_1}
%	\{
%		\norm{x - x_1}{0}^2
%		+
%		t^2
%		\norm{x_1}{1}^2
%	\}
%=
	\int_{[0, \infty)}
	\inf_{\epsilon > 0}
	\{
		(1 - \epsilon)^2
		+
		\epsilon^2 
		t^2 
		\lambda^{-2 \gamma}
	\}
	d\mu_{x}(\lambda)
=
	\int_{[0, \infty)}
	\frac{t^2}{t^2 + \lambda^{2 \gamma}}
	d\mu_{x}(\lambda)
	.
\end{align}
%
%
For the limiting cases $\nu = 0$ and $\nu = \gamma$
we recover from \eqref{e:0ga} and \eqref{e:KtI}
that
\begin{align}
\label{e:X01}
	X_{0 : \gamma} = X_0
	%
	\quad\text{and}\quad
	%
	X_{\gamma : \gamma} = X_\gamma
	%
	\quad
	\text{with equality of norms}
	.
\end{align}
%
%
%
%


% If $0 \leq \nu \leq \gamma$ then
% \begin{align}
% 	\normiii{x}{\nu} \leq \norm{x}{\gamma}
% \end{align}




The different spaces are related by the following Proposition.
% %
% The intermediate spaces
% $X_\gamma \subset X_{\nu:\gamma} \subset X$
% are characterized by the spectral measure:

\begin{proposition}
\label{p:ga-nu}
	Let $0 \leq \nu < \gamma$.
	%
% 	Set $\sigma := \norm{T^* T}{}$.
	%
	Let $x \in X$.
	%
	Then
	%
	\begin{align}
	\label{e:ga-nu}
		\sqrt{ 1 - \nu / \gamma }
		\norm{x}{\nu : \gamma}
	\stackrel{a)}{\leq}
		\normiii{x}{\nu}
	\stackrel{b)}{\leq}
		N_{\nu / \gamma}
		\norm{x}{\nu : \gamma}
	\stackrel{c)}{\leq}
		\norm{x}{\nu}
	\stackrel{d)}{\leq}
		\norm{T^* T}{}^{\gamma - \nu}
		\norm{x}{\gamma}
	.
	\end{align}
	%
	Hence, the following embeddings are continuous:
	\begin{align}
		X_{\gamma}
		\subset
		X_\nu
		\subset
		X_{\nu : \gamma}
		=
		\IX_\nu
		,
	\end{align}
	%
	where
	$X_{\nu : \gamma} = \IX_\nu$ 
	with equivalence of norms 
	possibly not uniform in $\nu < \gamma$.
\end{proposition}
%
%
\begin{proof}
	The inequality
	(\ref{e:ga-nu}d) follows
	from
	the representation \eqref{e:0ga},
	restricting the domain of integration
	to $[0, \norm{T^* T}{}]$.
	%
	The inequality (\ref{e:ga-nu}c)
	is obtained by 
	identifying $s := t \lambda^{-\gamma}$ and $\theta := \nu / \gamma$
	in \eqref{e:N=sup},
	and using it in \eqref{e:KtI},
	viz.~%
	\begin{align}
		\norm{x}{\nu : \gamma}^2 
		= 
		\sup_{t > 0} t^{-2 \nu/\gamma} |K_t^\gamma(x)|^2
		\leq
		N_{\nu / \gamma}^{-2}
		\int_{[0, \infty)} \lambda^{-2 \nu} d \mu_x(\lambda)
		.
	\end{align}
	%
	%
	%
	For the inequality (\ref{e:ga-nu}b)
	we use the identity
	\begin{align}
		s^{-2\nu}
		=
		N_{\nu/\gamma}^2
		t^{-2\nu/\gamma}
		\frac{ t^2 }{ t^2 + s^{2 \gamma} }
		\quad\text{for}\quad
		t = 
		s^\gamma
		\sqrt{ \frac{\gamma - \nu}{\nu} }
		,
	\end{align}
	combined with $\lambda \leq s$
	in the first step of 
	\begin{align*}
		s^{-2\nu}
		\norm{
			E_{[0, s)} x
		}{}^2
		%
		\leq
		%
		N_{\nu/\gamma}^2
		t^{-2\nu/\gamma}
		\int_{[0, s)}
		\frac{ t^2 }{ t^2 + \lambda^{2 \gamma} }
		d\mu_x(\lambda)
		%
		\stackrel{\eqref{e:KtI}}{\leq}
		%
		N_{\nu/\gamma}^2
		t^{-2\nu/\gamma}
		| K_t^\gamma(x) |^2
		.
	\end{align*}
	%
	Taking
	$\sup_{t > 0}$ on the right,
	then
	$\sup_{s > 0}$ on the left
	gives (\ref{e:ga-nu}b).
	%
	%
	%
	%
	Finally,
	from \eqref{e:KtI}, 
	followed by \eqref{e:Tail2},
	we have 
	\begin{align*}
		\norm{x}{\nu:\gamma}^2
	& \leq
		\sup_{t > 0} t^{-2 \nu/\gamma}
		\inf_{s > 0}
		\left\{
			\norm{
				E_{[0, s)} x
			}{}^2
			+
			s^2 \int_{[s, \infty)} \lambda^{-2\gamma} d \mu_x(\lambda)
		\right\}
	\leq
		\frac{\gamma}{\gamma - \nu}
		\normiii{x}{\nu}^2
		,
	\end{align*}
	with the choice $s := t^{1/\gamma}$ for the last inequality,
	and this shows (\ref{e:ga-nu}a).
	%
	%
	The last claim is a combination of (\ref{e:ga-nu}a) and (\ref{e:ga-nu}b).
\end{proof}

% The upper bound in \eqref{e:ga-nu} 
% may also be obtained more explicitly 

%

%
The choice $\gamma := 2 \nu$ in \eqref{e:ga-nu}
leads to the chain of inequalities
\begin{align}
\label{e:infg}
	\tfrac{1}{\sqrt{2}} \normiii{\cdot}{\nu}
	\leq
	\inf_{\gamma > \nu} \norm{\cdot}{\nu : \gamma}
	\leq
	\norm{\cdot}{\nu : 2 \nu}
	\leq
	\sqrt{2} \normiii{\cdot}{\nu}
	,
\end{align}
and therefore,
$X_{\nu : 2 \nu} = \IX_\nu$ 
with equivalence of norms uniformly in $\nu \geq 0$.
%
%
%
Equality between
$\inf_{\gamma > \nu} \norm{\cdot}{\nu : \gamma}$
and
$\norm{\cdot}{\nu : 2 \nu}$
% in \eqref{e:infg}
does not hold in general.
%
However, 
if $\mu_x$ is a Dirac measure supported at some $\lambda_0 > 0$, 
and $\nu > 0$,
then the infimum of
\begin{align}
	\norm{x}{\nu : \gamma}^2
	=
	\sup_{t > 0}
	t^{-2 \nu/\gamma}
	\frac{t^2}{t^2 + \lambda_0^{2 \gamma}}
	\stackrel{\eqref{e:N=sup}}{=}
	N_{\nu / \gamma}^{-2}
	\lambda_0^{-2 \nu}
	,
\end{align}
over $\gamma > \nu$
is indeed achieved at $\gamma = 2 \nu$.

%

The constants in (\ref{e:ga-nu}a) and (\ref{e:ga-nu}b) are sharp.
%
For example,
for
$\mu_x = \delta_{\lambda_0}$ being the Dirac measure at $\lambda_0 = 1$,
% % we find (using \eqref{e:N=sup})
\begin{align}
	\normiii{x}{\nu}
	=
	1
	%
	\quad\text{and}\quad
	%
	N_{\nu / \gamma}
	\norm{x}{\nu : \gamma}
	=
	1
	.
\end{align}
%
On the other hand,
for
$d \mu_x(\lambda) = 2 \nu \lambda^{2 \nu - 1} dx$
(that this measure is not compactly supported is not essential)
we find
$\normiii{x}{\nu} = 1$,
while
\begin{align}
	\norm{x}{\nu : \gamma}^2
	= 
	\sup_{t > 0} t^{-2\nu} |K_t^\gamma(x)|^2 = 
	\frac{\pi \nu / \gamma}{\sin(\pi \nu / \gamma)}
	=
	(1 + o(1))
	\frac{\gamma}{\gamma - \nu}
	\quad\text{as}\quad
	\nu \nearrow \gamma
	,
\end{align}
so that
% \begin{align}
% 	\normiii{x}{\nu}^2
% 	=
% 	1
% % 	=
% % 	\frac{
% % 		\sin(\pi \nu / \gamma)
% % 	}{
% % 		\pi \nu / \gamma
% % 	}
% % 	\norm{x}{\nu : \gamma}^2
% 	\norm{x}{\nu : \gamma}^2
% 	.
% \end{align}
% %
% This shows that 
the norm equivalence in \eqref{e:ga-nu}
does 
deteriorate
as $\nu \nearrow \gamma$.
%
The relation
to
the finite qualification of the (iterated) Tikhonov regularization
is discussed in
Section \ref{s:ip:T}.
%

%

We provide next another illustration
of the fact that $\IX_\nu$ is in general strictly larger than $X_\nu$,
here for $\nu := 1$.
%
We will revisit this example in Section \ref{s:ip:L}, Example \ref{x:IX:LW}.
%


\begin{example}
\label{x:IX}
	Consider the diagonal operator 
	$T := \mathop{\mathrm{diag}} ((n^{-1/2})_{n \geq 1})$ on
	the sequence space $X := \ell_2(\IN)$.
	%
	Then $T^* T$ has eigenvalues $\lambda_n = n^{-1}$, $n \geq 1$.
	%
	Let $x^\dag := (n^{-3/2})_{n \geq 1}$.
	%
	Then $\mu_{x^\dag} = \sum_{n \geq 1} n^{-3} \delta_{\lambda_n}$,
	so that
	\begin{align}
		\normiii{x^\dag}{1}^2 
		=
		\sup_{N \geq 1}
		\lambda_N^{-2} \mu_{x^\dag}([0, \lambda_N]) 
		= 
		\sup_{N \geq 1}
		N^2
		\sum_{n \geq N} n^{-3}
		< \infty
		,
	\end{align}
	yet,
	$\norm{x^\dag}{1}^2 = \sum_{n \geq 1} n^{-1}$ is not finite.
	%
	Therefore $x^\dag \in \IX_1 \setminus X_1$.
	%
	The $\sup$ is assumed at $N = 1$, 
	so that
	$\normiii{x^\dag}{1}^2 = \zeta(3) \approx (1.096)^2$.
\end{example}

%

\providecommand{\IX}{\mathbb{X}}
In \cite[Proposition 11]{AndreevElbauDeHoopLiuScherzer2015}
the variational inequality
\begin{align}
\label{e:S}
	\exists \beta \geq 0:
	\quad
	| \scalar{ x, \omega }{} |
	\leq
	\beta
	\norm{ (T^* T)^{\gamma} \omega }{}^{\nu / \gamma}
	\norm{ \omega }{}^{1 - \nu / \gamma}
	\quad
	\forall \omega \in X
	,
\end{align}
was shown 
for $0 < \nu < \gamma$
to hold 
if and only if
$x \in \IX_\nu$.
%
%
We prove a more precise statement,
in particular including
the limiting cases $\nu = 0$ and $\nu = \gamma$.
%
The first part of the proof
(the inequality ``$\leq$'')
simplifies and sharpens
the corresponding part of 
\cite[Proof of Proposition 11]{AndreevElbauDeHoopLiuScherzer2015}.
%
The second part (the inequality ``$\geq$'')
draws from \cite{Flemming2012}.
%

%

\begin{proposition}
\label{p:svi}
	Let $x \in X$ and $0 \leq \nu \leq \gamma$.
	Then
%	(note the asymmetry of subscripts),
	\begin{align}
	\label{e:p:svi}
		\sup_{\norm{\omega}{} = 1}
		\norm{(T^* T)^{\gamma} \omega}{}^{-\nu/\gamma}
		|\scalar{x, \omega}{}|
		=
		N_{\nu / \gamma}
		\norm{x}{\nu:\gamma}
	.
	\end{align}
	%
%	In particular, 
%	\eqref{e:S} holds
%	with $0 < \nu < \gamma$
%	if and only if $x \in \IX_\nu$.
%	,
%	and
%	a sufficient condition for \eqref{e:S} to hold
%	with $\nu = \gamma$
%	is
%	$x \in X_\gamma$.
\end{proposition}

\begin{proof}
	For $\nu = 0$ the statement is trivial
	due to
	$\norm{x}{0:\gamma} = \norm{x}{}$.
	%
	For $\nu = \gamma$
	the statement follows from \cite[Lemma 8.21]{ScherzerBook2009}
	and $\norm{x}{\gamma:\gamma} = \norm{x}{\gamma}$.
	%
	In both cases, recall $N_0 = N_1 = 1$.
	%
	For the remainder of the proof 
	we assume $0 < \nu < \gamma$.

	Let $\omega \in X$ 
	and
	consider the linear mapping $S : x \mapsto \scalar{ x, \omega }{}$.
	%
	Then
	$| S x | \leq \norm{\omega}{} \norm{x}{}$,
	and
	$| S x | \leq \norm{(T^* T)^{\gamma} \omega}{} \norm{x}{\gamma}$
	for all $x \in X_{\gamma}$.
	%
	By the operator interpolation inequality \eqref{e:|Snu|}
	we have
	$
		| \scalar{ x, \omega }{} |
	\leq
		N_{\nu/\gamma}
		\norm{x}{\nu : \gamma}
		\norm{(T^* T)^{\gamma} \omega}{}^{\nu/\gamma}
		\norm{\omega}{}^{1 - \nu/\gamma}
	$.
	%
	This implies ``$\leq$'' in \eqref{e:p:svi}.
	
	%
	
	To verify ``$\geq$'' in \eqref{e:p:svi},
	it suffices to establish the case $\gamma = 1$,
	then apply it 
	with $(T^* T)^\gamma$ replacing $T^* T$ 
	(also in the definitions of the norms in Section \ref{s:i:H}).
	%
	Thus we assume that \eqref{e:S} holds with $0 < \nu < \gamma = 1$.
	%
	To verify $N_{\nu} \norm{x}{\nu:1} \leq \beta$
	we check
	$D \leq N_{\nu}^{-2} \beta$
	for the quantity $D$ from \eqref{e:D=E}.
	%
	In other words
	we
	show the bound
	$d(r)^{1 - \nu} r^\nu \leq N_{\nu}^{-2} \beta$
	for
	the distance function $d(\cdot)$ from \eqref{e:dr}.
	%
	To that end we 
	combine \cite[Proposition 2.10]{Flemming2012},
	\cite[Theorem 4.1]{Flemming2012}
	and \cite[Theorem 4.5]{Flemming2012}:
	The inequality
	\begin{align}
		| \scalar{ x, \omega }{} |
		\leq
		\beta
		\norm{ (T^* T) \omega }{}^{\frac{\kappa}{2 - \kappa}}
		\quad
		\forall \omega \in X
		\quad\text{with}\quad
		\norm{\omega}{} = 1,
	\end{align}
	with 
	\begin{align}
		\kappa := 
		\tfrac{2 \nu}{1 + \nu}
		\quad\text{and}\quad
		\beta :=
		\tfrac{2 - \kappa}{2}
		\left(
			\tfrac{\tilde\beta}{1 - \kappa} 
		\right)^{\frac{1 - \kappa}{2 - \kappa}}
		a^{\frac{1}{2 - \kappa}}
		\quad
		\text{%
			for some $0 < \tilde\beta \leq 1$ and $a > 0$,
		}
	\end{align}
	implies
	$
		d^2(r) \leq 
		\tilde\beta \, (-\varphi)^*(-2 r)
	$
	for all $r > 0$,
	where $(-\varphi)^*$ is the Legendre--Fenchel conjugate
	of $t \mapsto -\varphi(t) := -a t^\kappa$ defined for $t > 0$.
	%
	Straightforward but tedious algebra
	yields
	the desired estimate 
	$d(r)^{1 - \nu} r^\nu \leq N_{\nu}^{-2} \beta$.
\end{proof}

%

%One can interpret the variational inequality \eqref{e:S}
%in yet another way.
%%
%Suppose $\gamma = 1$ for simplicity.
%%
%Noticing that
%any functional $f$ on $X_1$
%can be written as
%$f := \scalar{ \cdot, \omega }{}$
%with
%norm
%$\norm{f}{X_1'} = \norm{ (T^* T) \omega }{}$,
%we have from \eqref{e:S}
%the inequality
%\begin{align}
%	| f(x) |
%	\leq
%	\beta
%	\norm{ f }{X_1'}^{\nu}
%	\norm{ f }{X'}^{1 - \nu}
%	\quad
%	\forall f \in X_{1}'
%	.
%\end{align}

%\TODO{}
%The assumption $x \in X_{\nu : \gamma}$,
%or $x \in \IX_\nu$ combined with \eqref{e:ga-nu},
%now
%yields a quantitative version of \eqref{e:S}.
%%
%Thus, as $\nu \nearrow \gamma$,
%the equivalence 
%``\eqref{e:S} $\Leftrightarrow$ $x \in \IX_\nu$''
%deteriorates
%but 
%``\eqref{e:S} $\Leftrightarrow$ $x \in X_{\nu : \gamma}$''
%does not,
%and in fact remains valid for $\nu = \gamma$ by \cite[Lemma 8.21]{ScherzerBook2009}.
%%


% \TODO{Is $\inf_{\gamma > \nu} \norm{\cdot}{\nu : \gamma}$ a norm?}

%

%%%%%%%%%%%%%%%%%%%%%%%%%%%%%%%%%%%%%%%%%%%%%%%%%%%%%%%%%%%%%%%%%
%%%%%%%%%%%%%%%%%%%%%%%%%%%%%%%%%%%%%%%%%%%%%%%%%%%%%%%%%%%%%%%%%
\section{Application to linear inverse problems in Hilbert spaces}
\label{s:ip}
%%%%%%%%%%%%%%%%%%%%%%%%%%%%%%%%%%%%%%%%%%%%%%%%%%%%%%%%%%%%%%%%%
%%%%%%%%%%%%%%%%%%%%%%%%%%%%%%%%%%%%%%%%%%%%%%%%%%%%%%%%%%%%%%%%%

%%%%%%%%%%%%%%%%%%%%%%%%%%%%%%%%%%%%%%%%%%%%%%%%%%%%%%%%%%%%%%%%%
\subsection{Linear inverse problem}
\label{s:ip:0}
%%%%%%%%%%%%%%%%%%%%%%%%%%%%%%%%%%%%%%%%%%%%%%%%%%%%%%%%%%%%%%%%%

Let $X$ and $Y$ be Hilbert spaces.
%
Let $T : X \to Y$ be a bounded linear operator,
with possibly nonclosed range.
%
We write $\norm{\cdot}{}$ for the norm of $X$ and for that of $Y$.
%
As in Section \ref{s:i:H},
we assume that $T$ is injective.
%
Fix $y^\dag \in T X$ and 
let $x^\dag$ denote the solution to
\begin{align}
\label{e:Tx=y}
	T x^\dag = y^\dag
	.
\end{align}
%
% for which $\norm{x^\dag - x_0}{}$ is minimal,
Let
$ %\begin{align}
	x_0 \in X,
$ %\end{align}
called a prior,
be given.
%
%
%
The task is to find an approximation of $x^\dag$,
given $y^\dag$ (noise-free case)
or $y^\delta \approx y^\dag$ (noisy case)
with
\begin{align}
\label{e:ydel}
	\norm{ y^\dag - y^\delta }{}
	\leq
	\delta
	.
\end{align}


% Norm-minimality of $x^\dag - x_0$
% allows (\TODO{really?})
% to replace $X$ by $X / \ker T$ with the usual quotient norm
% so that $T$ 
% may be assumed to be injective.
%
We use the notation
from Section \ref{s:i:H},
including
the spectral measure $E$, 
the Borel measure $\mu_x$, 
the spaces $X_\gamma$, $X_{\nu:\gamma}$, etc.
%


%%%%%%%%%%%%%%%%%%%%%%%%%%%%%%%%%%%%%%%%%%%%%%%%%%%%%%%%%%%%%%%%%
\subsection{Spectral cut-off regularization}
\label{s:ip:cut}
%%%%%%%%%%%%%%%%%%%%%%%%%%%%%%%%%%%%%%%%%%%%%%%%%%%%%%%%%%%%%%%%%

The spectral cut-off regularization of $x^\dag$ is defined
as $x_0 + E_{[\alpha, \infty)} (x^\dag - x_0)$
for a parameter $\alpha > 0$.
%
From the definition \eqref{e:normiii} of $\normiii{\cdot}{\nu}$
it is immediate that
the error of this regularization is 
\begin{align}
\label{e:coerr}
	\norm{(I - E_{[\alpha, \infty)}) (x^\dag - x_0)}{}
	=
	\norm{E_{[0, \alpha)} (x^\dag - x_0)}{} 
	\leq 
	\alpha^\nu \normiii{x^\dag - x_0}{\nu}
	\quad
	\forall \alpha \geq 0
	\quad
	\forall \nu \geq 0
	.
\end{align}
%
Since there are no restrictions on the possible convergence rate $\nu \geq 0$
(referred to as infinite qualification),
and no further constants are involved,
we may view the performance of this regularization as a reference.

%%%%%%%%%%%%%%%%%%%%%%%%%%%%%%%%%%%%%%%%%%%%%%%%%%%%%%%%%%%%%%%%%
\subsection{Tikhonov regularization}
\label{s:ip:T}
%%%%%%%%%%%%%%%%%%%%%%%%%%%%%%%%%%%%%%%%%%%%%%%%%%%%%%%%%%%%%%%%%

For $\alpha > 0$, 
the regularized solution $x_\alpha \in X$ in the noise-free case
is defined 
as the unique minimizer of the Tikhonov functional
\begin{align}
\label{e:J}
	J_\alpha(x; x_0)
	:=
	\norm{ y^\dag - T x }{}^2
	+
	\alpha
	\norm{ x - x_0 }{}^2,
	\quad
	x \in X
	.
\end{align}
%
Replacing $y^\dag$ by $y^\delta$
defines
the regularized solution $x_\alpha^\delta \in X$ 
in the noisy case.
%
They are equivalently characterized by
the first order optimality conditions
\begin{align}
\label{e:xx}
	x_\alpha = (T^* T + \alpha I)^{-1} (T^* y^\dag + \alpha x_0)
	\quad\text{and}\quad
	x_\alpha^\delta = (T^* T + \alpha I)^{-1} (T^* y^\delta + \alpha x_0)
	.
\end{align}


Writing $e^\delta := x_\alpha - x_\alpha^\delta$ for the moment,
we have
$ %\begin{align*}
	\norm{ T e^\delta }{}^2 + \alpha \norm{ e^\delta }{}^2
=
	\scalar{ (T^* T + \alpha I) e^\delta, e^\delta }{}
=
	\scalar{ T^* (y^\dag - y^\delta), e^\delta }{}
\leq
	\norm{ y^\dag - y^\delta }{} \norm{ T e^\delta }{}
\leq
	\tfrac14 \delta^2 + \norm{ T e^\delta }{}^2
	,
$ %\end{align*}
and cancellation of $\norm{ T e^\delta }{}^2$ on both ends gives
the error splitting 
\begin{align}
\label{e:xad}
	\norm{ x^\dag - x_\alpha^\delta }{}
	\leq
	\norm{ x^\dag - x_\alpha }{}
	+
	\norm{ e^\delta }{}
	\leq
	\norm{ x^\dag - x_\alpha }{}
	+
	\tfrac12
	\tfrac{\delta}{\sqrt{\alpha}}
	.
\end{align}
%
%Note that \eqref{e:xad} is independent of 
%the prior $x_0$.

The parameter $\alpha > 0$ is determined by
a parameter choice strategy
\begin{align}
\label{e:pcs}
	\bar{\alpha} : 
	(\delta, y^\delta, \ldots) \mapsto \bar{\alpha}(\delta, y^\delta, \ldots)
	.
\end{align}
%
The one that minimizes 
$\alpha \mapsto \norm{x^\dag - x_\alpha^\delta}{}$
whenever $y^\delta$ and $y^\dag$ are fixed
may be considered the optimal strategy.
%
We shall suppose that
the parameter choice strategy 
satisfies
\begin{align}
\label{e:qo}
%	\sup_{\delta > 0}
	\delta^{-2\nu}
	\sup_{{{\eqref{e:ydel}}}}
	\norm{x^\dag - x_{\bar{\alpha}(\delta, y^\delta, \ldots)}^\delta}{}^{2\nu + 1}
	\sim
	\sup_{\alpha > 0}
	\alpha^{-\nu}
	\norm{x^\dag - x_{\alpha}}{}
	\quad
	\forall
	\delta > 0
	,
\end{align}
where the hidden constants 
do not depend on $\delta$, the exact solution $x^\dag$, or the prior $x_0$,
but may depend on $\nu \geq 0$.
%
Here, $\sup_{{{\eqref{e:ydel}}}}$
means the supremum over all $y^\delta \in Y$
which satisfy \eqref{e:ydel}.
%
As an example,
the a priori parameter choice strategy
$\bar{\alpha}$
% only depending on $\delta$ and $y^\dag$, and
defined by
\begin{align}
\label{e:pcs1}
	\sqrt{\bar{\alpha}(\delta, y^\dag)}
	\norm{ x^\dag - x_{\bar{\alpha}(\delta, y^\dag)} }{}
	\stackrel{{!}}{=}
	\tfrac12
	\delta
	\quad
	\forall \delta > 0
\end{align}
satisfies \eqref{e:qo}.
%
Specifically, 
the error splitting
\eqref{e:xad} quickly yields
$\text{LHS} \leq 2 \, \text{RHS}$ in \eqref{e:qo}
and
an inspection of
\cite[Proof of Theorem 2.6]{Neubauer1997}
yields
$\frac{13}{2} (5 / \sqrt{2})^{2 \nu + 1} \text{LHS} \geq \text{RHS}$.
%
That proof
also shows that the optimal strategy, see above,
satisfies \eqref{e:qo}.
%
Of course, of practical interest are 
parameter choice strategies
that do not access 
the exact data $y^\dag$ or the exact solution $x^\dag$;
\REV{%
in that regard 
the notion 
of quasioptimality by Raus \& H\"amarik \cite{RausHamarik2007}
is useful,
see comments following Proposition \ref{p:tik} below.
}%\REV
% such as \cite{Gfrerer1987}.
%
% \TODO{examples; q-o?}
% \TODO{see Hegland, ``An optimal ...'', 1992 for refs}
%
In any case, 
\eqref{e:qo} formalizes
the equivalence
\begin{align}
	\text{``}
		\norm{ x^\dag - x_{\bar{\alpha}(\delta, y^\delta, \ldots)}^\delta }{}
		\lesssim
		\delta^{2 \nu / (2 \nu + 1)}
		\quad
		\forall \delta > 0
	\text{''}
	\quad\Leftrightarrow\quad
	\text{``}
		\norm{x^\dag - x_\alpha}{}
		\lesssim
		\alpha^\nu
		\quad
		\forall \alpha > 0
	\text{''}
\end{align}
of the error estimates 
in the noisy and in the noise-free 
cases,
and
in the following we assume \eqref{e:qo},
and focus on its \text{RHS}.
%
%Note that \eqref{e:qo} is 
%invariant under the choice of
%the prior $x_0$,
%which is of importance 
%in
%the iterated Tikhonov regularization discussed below.

%

From the abstract theory of interpolation,
convergence rates of 
the Tikhonov regularization error $\norm{x^\dag - x_\alpha}{}$
for $x^\dag \in X_{\nu:1}$ when $0 < \nu < 1$
quickly follow.
%
Indeed,
for $\alpha > 0$
consider the linear mapping
$S_\alpha : X \to X$,
$(x^\dag - x_0) \mapsto (x^\dag - x_\alpha)$.
%
From \eqref{e:xx},
\begin{align}
\label{e:Sa}
	S_\alpha = \alpha (T^* T + \alpha I)^{-1}
	,
\end{align}
so that
$\norm{S_\alpha }{} \leq 1$.
%
%Suppose $x_0 = 0$ for simplicity.
%
Under the classical source condition
\begin{align}
\label{e:cla}
	x^\dag \in 
	x_0 + X_1
	\quad\text{where}\quad
	X_1 = 
	\mathop{\mathrm{range}} (T^* T)
	,
\end{align}
it is known
(and shown below in \eqref{e:p:tik} for $\nu = 1$)
that
\begin{align}
\label{e:a01}
	\norm{ S_\alpha (x^\dag - x_0) }{} 
	\leq 
	\alpha \norm{x^\dag - x_0}{1}
	.
\end{align}
%
The operator interpolation inequality 
\eqref{e:|Snu|} implies
\begin{align}
\label{e:anu}
	\norm{ x^\dag - x_\alpha }{} 
	=
	\norm{ S_\alpha (x^\dag - x_0) }{}
	\leq
	\alpha^\nu
	N_\nu
	\norm{ x^\dag - x_0 }{\nu:1}
	\quad
	\forall x^\dag \in x_0 + X_{\nu:1}
	\quad
	\forall \alpha > 0
	.
\end{align}
%
The following Proposition
%refinement of this statement
% \ref{p:tik} below
shows that $x_0 + X_{\nu:1} \subset X$ is precisely
the (affine) subspace that allows those convergence rates,
which is the main observation of this note.


\begin{proposition}
\label{p:tik}
	Let $0 \leq \nu \leq 1$. Then
	\begin{align}
	\label{e:p:tik}
%		\tfrac{1}{\sqrt{2}}
		N_\nu^{-1}
		\norm{x^\dag - x_0}{\nu:1}
		\leq
		\sup_{\alpha > 0}
		\alpha^{-\nu}
		\norm{ x^\dag - x_\alpha }{}
		\leq
		\norm{x^\dag - x_0}{\nu:1}
		.
	\end{align}
	%
	Equalities hold for $\nu = 0$ and $\nu = 1$.
	%
	If $\nu = 1$ then $\sup_{\alpha > 0}$ can be replaced by $\lim_{\alpha \searrow 0}$.
\end{proposition}

%
The result is a special case of Proposition \ref{p:tikk} below,
and the proof is therefore omitted.
%
Several remarks are in order.

%\begin{enumerate}
%\item
	
	Combining (\ref{e:ga-nu}a)--(\ref{e:ga-nu}b) and \eqref{e:p:tik},
	for $0 < \nu < 1$,
	we have
	that $\norm{x^\dag - x_\alpha}{} \leq C \alpha^{\nu}$ for all $\alpha > 0$
	if and only if
	$x^\dag \in x_0 + \IX_\nu$.
	%
	In essence,
	this was already shown in \cite[Theorem 2.1]{Neubauer1997}.
	%
	%
	We emphasize, however,
	that (\ref{e:ga-nu}a) and \eqref{e:p:tik} yield
	the more precise upper bound
	\begin{align}
	\label{e:bc}
		\norm{x^\dag - x_\alpha}{} 
		\leq 
		\alpha^\nu \tfrac{1}{\sqrt{1 - \nu}} 
		\normiii{x^\dag - x_0}{\nu}
		\quad
		\forall \alpha > 0 
		.
	\end{align}
	%
	In particular, 
	although the rate of convergence of Tikhonov regularization
	is at least $\nu$ whenever $x^\dag \in x_0 + \IX_\nu$,
	the constant in \eqref{e:bc} may deteriorate as $\nu \nearrow 1$
	compared to 
	the error \eqref{e:coerr}
	of the spectral cut-off regularization.
	%
	This may be interpreted as a quantitative description of 
	the finite qualification of Tikhonov regularization,
	that is its inability to provide convergence larger than $\nu = 1$.
	%
	Of course,
	by \eqref{e:p:tik},
	the rate of $\nu = 1$ does hold 
	if (and only if) $x^\dag \in x_0 + X_1$,
	but
	this condition is more restrictive 
	than $x^\dag \in x_0 + \IX_1$,
	cf.~%
	Proposition \ref{p:ga-nu}
	and
	Example \ref{x:IX}.
	
	%
	

%\item

	\REV{%
	From \eqref{e:p:tik}
	we see that
	Tikhonov regularization with 
	a parameter choice rule satisfying \eqref{e:qo}
	is order optimal
	on 
	$M_{\nu,\rho} := \{ x \in X : \norm{x}{\nu:1} \leq \rho \}$
	if $T$ has nonclosed range.
	%
	Indeed, \eqref{e:qo} and \eqref{e:p:tik}
	imply
	the error estimate
	$
		\norm{x^\dagger - x_{\bar{\alpha}}^\delta }{}
		\lesssim
		\delta^{2\nu / (2\nu + 1)}
		\rho^{1 / (2\nu + 1)}
	$
	whenever
	$x^\dagger \in M_{\nu,\rho}$;
	on the other hand,
	\cite[Prop 3.15]{EnglHankeNeubauer1996}
	and comments there
	show that
	this estimate is optimal
	because
	$M_{\nu,\rho}$ contains
	the classical source set
	$\{ x \in X : \norm{x}{\nu} \leq \rho' \}$
	by (\ref{e:ga-nu}c).
	%
	%
	Now we can invoke
	\cite[Thm 2.3]{RausHamarik2007}
	to assert that
	any parameter choice rule 
	that is strongly quasioptimal
	in the sense of \cite[Def 2.2]{RausHamarik2007}
	will again give rise 
	to
	an order optimal method on $M_{\nu,\rho}$.
	}%\REV

%\item

	The final statement of Proposition \ref{p:tik}
	implies the saturation result:
	If $\norm{x^\dagger - x_\alpha}{} = o(\alpha)$
	as $\alpha \searrow 0$
	then $x^\dagger = x_0$.
	%
	An analogous statement holds 
	in the noisy case 
	for any
	parameter choice strategy
	satisfying \eqref{e:qo}.
	
%\item

	An additional consequence 
	of \eqref{e:p:tik}
	is that,
	for $0 \leq \nu \leq \gamma$,
	\begin{align}
	\label{e:c:tikr}
		N_{\nu/\gamma}^{-1}
		\norm{ x^\dag - x_0 }{\nu : \gamma}
		\leq
		\sup_{\alpha > 0}
		\alpha^{-\nu}
		\norm{ x^\dag - x_{\alpha:\gamma} }{}
		\leq
		\norm{ x^\dag - x_0 }{\nu : \gamma}
	\end{align}
	where $x_{\alpha:\gamma}$ is given by
	\begin{align}
	\label{e:c:tikr:x}
		x_{\alpha:\gamma} 
		:= 
		\mathop{\mathrm{argmin}}_{x \in X} 
		\{
			\norm{ (T^* T)^{\gamma/2} (x^\dag - x) }{}^2
			+
			\alpha
			\norm{ x^\dag - x_0 }{}^2
		\}
		.
	\end{align}
	%
	This is simply \eqref{e:p:tik}
	for the operator
	$(T^* T)^{\gamma/2}$ instead of $T$.
	%
	%
	%
	Therefore,
	letting $x_0 = 0$,
	the mapping
	$x^\dag \mapsto \sup_{\alpha > 0} \alpha^{-\nu} \norm{x^\dag - x_{\alpha:\gamma}}{}$
	defines a norm on $X_{\nu:\gamma}$
	which 
%	in view of \eqref{e:c:tikr}
	is equivalent to $\norm{\cdot}{\nu:\gamma}$
	uniformly in $0 \leq \nu \leq \gamma$.
	%

% \item

	The variational inequality \eqref{e:S}
	was used as a starting point in 
	\cite[Lemma 2(ii) \& Lemma 7]{AndreevElbauDeHoopLiuScherzer2015}
	for an elementary proof of
	convergence rates for Tikhonov regularization 
	without 
	invoking spectral theory.
	%
	This is certainly of interest,
	but leads to the natural question 
	whether \eqref{e:S} itself
	can be established in particular cases 
	without spectral theory; 
	the characterization \eqref{e:p:svi}
	in terms of interpolation spaces
	% given in Proposition \ref{p:svi}
	is a step in that direction.
	%
	That elementary proof, however,
	seems to be limited to rates $\nu \leq 1/2$ 
	in \eqref{e:anu},
	as it uses the variational inequality \eqref{e:S} 
	with $\gamma = 1/2$.

%\item
	
	As an example
	consider
	the following integral operator
	on
	% average-free 
	square integrable functions
	$X := Y := L_2$
	of the unit interval $(-1, 1)$,
	%suppressed in the notation,
% 	into $Y := L_2$,
	\begin{align}
		\textstyle
		T : X \to Y,
	% 	\quad
	% 	x \mapsto T x,w2e3
		\quad
		(T x)(t)
		:=
		\int_{-1}^{t} x(s) ds
		-
		\frac{1}{2}
		\int_{-1}^1 
			\int_{-1}^\tau x(s) ds
		d\tau
		.
	\end{align}
	%
	One can check that
% 	\begin{align}
% 		\textstyle
% 		(T^* T x)(s)
% 		=
% 		\int_s^1 \int_{-1}^t x(\sigma) d\sigma dt
% 		+
% 		\frac{s-1}{2}
% 		\int_{-1}^1 \int_{-1}^t x(\sigma) d\sigma dt
% 		,
% 	\end{align}
	$u := T^* T x$
	solves
	the boundary value problem
	$-u'' = x$, $u|_{\pm1} = 0$.
	%
	%
	Hence, 
	$X_1$ from \eqref{e:Xga}
	is 
	the Sobolev space $H^2 \cap H_0^1$.
% 	It is equipped with the norm $\norm{(\cdot)''}{L_2}$.
	%
	%
	Let now $y^\dagger(t) := |t| - 1$,
	so that \eqref{e:Tx=y}
	holds
	with $x^\dagger(s) := \mathop{\mathrm{sign}}(s)$.
	%
	This means, 
	solving \eqref{e:Tx=y} with noisy data $y^\delta$
	amounts to taking the derivative of $y^\dagger$
	perturbed by $L_2$ noise.
	%
	Set $\nu := 1/4$.
	%
	The interpolation space
	$X_{\nu:1} = (X, X_1)_{1/4, \infty}$
	is the Besov space $\mathring{B}_{2, \infty}^{1/2}$
	\cite[Prop 1.25 and 1.30]{Guidetti1991}.
	%
	In view of \cite[Def 1.12 of $\mathring{B}$]{Guidetti1991},
	%
	to check that $x^\dagger \in X_{\nu:1}$,
	it suffices to check that
	$L_2$ step functions are in 
	$B_{2,\infty}^{1/2}(\IR)$,
	but
	this is immediate from
	the intrinsic definition 
	\cite[Def 1.1]{Guidetti1991}.
	%
	On the other hand,
	\cite[Ch 1, Thm 11.7]{LionsMagenes1972}
	implies that
	$t \mapsto |x(t)|^2 / (t^2 - 1)$ is integrable
	for any $x \in X_{1/4}$
	(denoted by $H_{00}^{1/2}$ in \cite{LionsMagenes1972}),
	which is clearly not the case for $x^\dagger$.
	%
	In summary,
	$x^\dagger \in X_{(1/4):1} \setminus X_{1/4}$,
	and 
	Proposition \ref{p:tik}
	implies 
	the  
	convergence rate $\nu = 1/4$,
	familiar from \cite[Example 5]{FlemmingHofmannMathe2011}.

\REV{%
	On a more general note,
	let $M$ be
	a connected smooth finite-dimensional Riemannian manifold
	with bounded geometry
	(complete, injectivity radius bounded away from zero, bounded covariant derivatives of the curvature tensor). 
	%
	For instance, 
	the flat space $\IR^d$
	or
	any compact manifold without boundary
	has bounded geometry.
	%
	On $M$,
	Sobolev $H^s$ and Besov $B_{2,\infty}^s$
	spaces can be defined 
	intrinsically through local means \cite[\S1.11.4]{Triebel1992},
	and by \cite[\S1.11.3]{Triebel1992}
	one has
	$(H^s, H^{s+k})_{\theta, \infty} = B_{2,\infty}^{s + \theta k}$.
	%
	Thus, 
	for $T^* T$ with domain $X := H^s$ and range $X_1 = H^{s + 2m}$,
	$m > 0$,
	we have
	$X_{\nu:1} = B_{2,\infty}^{s + \nu 2 m}$,
	so that the convergence rate $\nu$ in \eqref{e:anu}
	is characterized in terms of Besov smoothness of $x^\dagger$.
	%
	This covers a wide class of operators $T$,
	for instance 
	bijective elliptic pseudo-differential operators
	$H^s \to H^{s+m}$
	of order $-m$
	such as
	$(I - \Delta)^{-m/2}$.
}%\REV
	
%\end{enumerate}


%%%%%%%%%%%%%%%%%%%%%%%%%%%%%%%%%%%%%%%%%%%%%%%%%%%%%%%%%%%%%%%%%
\subsection{Stationary iterated Tikhonov regularization}
\label{s:ip:IT}
%%%%%%%%%%%%%%%%%%%%%%%%%%%%%%%%%%%%%%%%%%%%%%%%%%%%%%%%%%%%%%%%%

Fixing $\alpha > 0$,
for each integer $k \geq 1$,
define $x_{\alpha, k}$ as the minimizer of
the Tikhonov functional
$x \mapsto J_\alpha(x; x_{\alpha,k-1})$,
where $J_\alpha$ is given by \eqref{e:J}
and
$x_{\alpha,0} := x_0$.
%
Thus, each iterate is used as a prior 
for the next one
with the same (stationary) choice of 
the regularization parameter $\alpha > 0$.
%
%
%
Then 
\begin{align}
\label{e:errTk}
	x^\dagger - x_{\alpha,k} = S_\alpha^k (x^\dagger - x_0)
	\quad
	\forall k \geq 1,
\end{align}
with $S_\alpha$ from \eqref{e:Sa}.
%
%
%Since, for any integer $k \geq 1$,
%the inequalities
%\begin{align}
%	2^{1 - 2k}
%	\frac{1}{1 + r^{2k}}
%	\leq
%	\frac{1}{(1 + r)^{2k}} 
%	\leq
%	\frac{1}{1 + r^{2k}}
%	\quad
%	\forall r \geq 0
%\end{align}
%hold (applied to $r = \lambda / t^{1/k}$),
%
%
%
Since $S_\alpha$ commutes with powers of $(T^* T)$,
we obtain from \eqref{e:a01}
the ``shift estimate''
$
	\norm{ S_\alpha (x^\dag - x_\alpha) }{k - \ell} 
	\leq 
	\alpha \norm{x^\dag - x_0}{k - \ell + 1}
$,
and repeating this argument for each iteration $\ell = 1, 2, \ldots, k$,
of the $k$-fold Tikhonov regularization
then
$ %\begin{align}
% \label{e:a0k}
	\norm{ S_\alpha (x^\dag - x_{\alpha, k}) }{} 
	\leq
	\alpha^k
	\norm{x^\dag - x_0}{k}
	.
$ %\end{align}
%
%
The operator interpolation inequality 
readily yields
the convergence rate
$\alpha^{\nu}$ 
for $x^\dag \in x_0 + X_{\nu:k}$ and $0 \leq \nu \leq k$.
%
The following extension of 
Proposition \ref{p:tik}
is more precise.

\begin{proposition}
\label{p:tikk}
	Let $k \geq 1$ be an integer, let $0 \leq \nu \leq k$ be real.
	%
	Then
	\begin{align}
	\label{e:p:tikk}
%		\sqrt{2}^{1 - 2 k}
		N_{\nu / k}^{1 - 2 k}
		\norm{ x^\dag - x_0 }{\nu : k}
		\leq
		\sup_{\alpha > 0}
		\alpha^{-\nu}
		\norm{ x^\dag - x_{\alpha,k} }{}
		\leq
		\norm{ x^\dag - x_0 }{\nu : k}
		.
	\end{align}
	%
% 	Equalities hold for $\nu = 0$ and $\nu = k$.
	%
	If $\nu = k$ then $\sup_{\alpha > 0}$ can be replaced by $\lim_{\alpha \searrow 0}$.
\end{proposition}
%
%
\begin{proof}
	From \eqref{e:errTk} and \eqref{e:Sa},
	\begin{align}
	\label{e:p:tikk:a}
		\alpha^{-2 \nu}
		\norm{ x^\dag - x_{\alpha,k} }{}^2
	 	=
	 	\alpha^{-2 \nu}
	 	\int_{[0, \infty)}
	 	\left( \frac{\alpha}{\alpha + \lambda} \right)^{2 k}
	 	d \mu_{x^\dag - x_0}(\lambda)
		.
	\end{align}
	%
	The first inequality in \eqref{e:p:tikk}
	is therefore a consequence of \eqref{e:Ncomm}
	with $\theta = \nu / k$,
	and the representation \eqref{e:KtI} of the interpolation norm.
	%
	The second inequality in \eqref{e:p:tikk}
	follows by identifying $\alpha = t^{1/k}$,
	estimating
	$1 / (\alpha + \lambda)^{2 k} \leq 1 / ( t^2 + \lambda^{2 k} )$,
	and
	invoking the representation \eqref{e:KtI}.
	%
	Finally,
	if $\nu = k$,
	then
	$\sup_{\alpha > 0}$\eqref{e:p:tikk:a} = 
	$\lim_{\alpha \searrow 0}$\eqref{e:p:tikk:a}
	by
	the Lebesgue monotone convergence theorem.
\end{proof}


%Since \eqref{e:qo}
%is invariant under the choice of the prior,
%we may replace 
%the couple
%$\{ x_{\bar{\alpha}(\delta)}, x_{\bar{\alpha}(\delta)}^\delta \}$
%therein by
%the couple
%$\{ x_{\bar{\alpha}(\delta),k}, x_{\bar{\alpha}(\delta),k}^\delta \}$.
%%
%For example,
%for 
%the
%parameter choice strategy \eqref{e:pcs1},
%we obtain
%from \eqref{e:p:tikk},
%for any $0 \leq \nu \leq k$,
%the estimate
%\begin{align}
%	\norm{ x^\dag - x_{\bar{\alpha}(\delta),k}^\delta }{} 
%	\leq 
%	\delta^{2 \nu / (2 \nu + 1)}
%	(
%		2
%		\norm{ x^\dag - x_0 }{\nu : k}
%	)^{1 / (2 \nu + 1)}
%	\quad
%	\forall 
%	\delta > 0
%\end{align}
%for the regularization error 
%in the noisy case.

%

Remarks analogous to those following Proposition \ref{p:tik} apply.
%
For instance,
the $k$-fold Tikhonov regularization
saturates:
If $\norm{x^\dag - x_{\alpha, k}}{} = o(\alpha^k)$ as $\alpha \searrow 0$
then $x^\dag = x_0$.

%Combining \eqref{e:p:tikk} with \eqref{e:ga-nu} we have
%for $0 < \nu < k$,
%\begin{align}
%	\norm{ x^\dag - x_{\alpha,k} }{}
%	\leq
%	\alpha^\nu
%	\sqrt{ \tfrac{k}{k - \nu} }
%	\normiii{x^\dag - x_0}{\nu}
%	.
%\end{align}
%
%The constant deteriorates as $\nu \nearrow k$,
%see discussion following Proposition \ref{p:tik}.



%%%%%%%%%%%%%%%%%%%%%%%%%%%%%%%%%%%%%%%%%%%%%%%%%%%%%%%%%%%%%%%%%
\subsection{Landweber iteration}
\label{s:ip:L}
%%%%%%%%%%%%%%%%%%%%%%%%%%%%%%%%%%%%%%%%%%%%%%%%%%%%%%%%%%%%%%%%%

Fix $\sigma > 0$ with
\begin{align}
\label{e:tau}
	0 < \sigma \norm{T^* T}{} \leq 1
	.
\end{align}
%
The Landweber iterates $x_k$ are defined by
\begin{align}
\label{e:LW}
	x_k :=
	x_{k-1} + \sigma T^* (y^\dag - T x_{k-1}),
	\quad
	k \in \IN,
	\quad\text{with a given}\quad
	x_0 \in X,
\end{align}
and $x_k^\delta$ by the same iteration with $y^\dag$ replaced by $y^\delta$
in the noisy case.
%
%
%
%
%The iteration \eqref{e:LW} is equivalent
%to the same iteration without $\sigma$,
%applied to 
%the equation $(\sqrt{\sigma} T) (\tfrac{1}{\sqrt{\sigma}} x) = y^\dag$.
%
%
One motivation for this iteration
is that
$x^\dag$ is a fixed point.
%
% thus, 
% subtracting \eqref{e:LW} from the corresponding equation with $x^\dag$,
By induction
one finds
\begin{align}
	x^\dag - x_k
	=
	(I - \sigma T^* T)
	(x^\dag - x_{k-1})
	=
	(I - \sigma T^* T)^k 
	(x^\dag - x_0)
	\quad
	\forall
	k \in \IN
	,
\end{align}
and similarly
the residual representation
\begin{align}
\label{e:dade}
	y^\dag - T x_k
	=
	(I - T T^*)^k (y^\dag - T x_0)
	%
	\quad\text{and}\quad
	%
	y^\delta - T x_k^\delta
	=
	(I - T T^*)^k (y^\delta - T x_0)
	.
\end{align}
%
%
The condition \eqref{e:tau} 
therefore guarantees
nondivergence of the iterates.

For the noisy case, 
an error splitting analogous to \eqref{e:xad}
is true \cite[Lemma 6.2]{EnglHankeNeubauer1996}:
\begin{align}
\label{e:LW:es}
	\norm{ x^\dag - x_k^\delta }{}
	\leq
	\norm{ x^\dag - x_k }{}
	+
	\delta
	\sqrt{k}
	\quad
	\forall k \geq 0
	,
\end{align}
thus one often ``morally'' identifies 
$k$ with $1/\alpha$.
%
%
We call a mapping 
\begin{align}
\label{e:sr}
	\bar{k}
	:
	(\delta, y^\delta, \ldots)
	\mapsto
	\bar{k}
	(\delta, y^\delta, \ldots)
\end{align}
a stopping rule,
and
in analogy to \eqref{e:pcs}--\eqref{e:qo}
we shall suppose that
it satisfies
\begin{align}
\label{e:qok}
	\delta^{-2\nu}
	\sup_{{{\eqref{e:ydel}}}}
	\norm{x^\dag - x_{\bar{k}(\delta, y^\delta, \ldots)}^\delta}{}^{2\nu + 1}
	\sim
	\sup_{k \geq 0}
	(1 + k/\nu)^{\nu}
	\norm{x^\dag - x_k}{}
	\quad
	\forall
	\delta > 0
	,
\end{align}
where the hidden constants 
do not depend on $\delta$, $x^\dag$, or $x_0$,
but may depend on $\nu \geq 0$.
%
The factor $(1 + k/\nu)$ 
is motivated by \eqref{e:LWE}--\eqref{e:p:LW}.
%
The stopping rule $\bar{k}$ may be based 
on the knowledge of some of the iterates $x_k$,
for example 
it may be the smallest $k \geq 0$ for which
the 
% Morozov 
discrepancy principle
\begin{align}
\label{e:MDPk}
	\norm{
		y^\delta - T x_k^\delta
	}{}
	\leq
	\delta \tau
\end{align}
is satisfied with some fixed threshold $\tau > 1$,
which in particular is not allowed to depend on $x^\dag$.
%
We will comment on this stopping rule at the end of this section.
%
%
% Again, we assume henceforth 
% that \eqref{e:qok} holds
% in order to focus on the noise-free case.



%
We shall show that 
convergence rates of 
the Landweber iteration \eqref{e:LW}
in the noise-free case
characterize the spaces $\IX_\nu$.
%
Before formalizing this, we need a lemma.

\begin{lemma}
\label{l:c2}
	The constant 
	\begin{align}
	\label{e:c2}
		c_2
		:=
		\sup_{a > 0}
		\sup_{s > 0}
		\sqrt{ I(s, a) },
		\quad
		I(s, a)
		:=
		\left(
			\frac{ s + a }{ a }
		\right)^a
		\int_0^1 s (1 - t)^{s-1} t^{a} dt,
	\end{align}
	is finite
	with
	$c_2 \approx 1.135$.
\end{lemma}
%
\begin{proof}
	First, $I(s, a)$ is well-defined
	for all real $s > 0$ and $a > 0$.
	%
	The integral, 
	known as the beta function \cite[\S6.2]{AbramowitzStegun1972},
	evaluates to 
	$\Gamma(s+1) \Gamma(a+1) / \Gamma(s+a+1)$.
	%
	One has $\lim_{s \searrow 0} I(s, a) = 1$ for all $a > 0$.
	%
	Moreover,
	the function
	$s \mapsto I(s, a)$
	is increasing for $0 \leq a \leq 1$
	and
	decreasing for $1 \leq a$.
	%
	So we need to consider 
	the pointwise limit $I(s, a)$ as $s \to \infty$
	for $0 \leq a \leq 1$.
	%
	Stirling's formula \cite[\S6.1.38]{AbramowitzStegun1972} 
	for $\Gamma(\cdot)$ yields
	\begin{align}
		\label{e:soo}
		\bar{I}(a) := \lim_{s \to \infty} I(s, a) = a^{-a} \Gamma(a + 1)
		\quad
		\forall a > 0
		,
	\end{align}
	%
	which clearly is bounded in $0 \leq a \leq 1$.
	%
	We find
	numerically
%	\footnote{%
%		For example, type
%%		``\texttt{lim ((s+a)/a){\^{}}a gamma(a+1) gamma(s+1) / gamma(s+a+1) as s to inf}''
%%		and
%		``\texttt{max gamma(1+x)/x{\^{}}x for x > 0}''
%		in \url{http://www.wolframalpha.com/}.
%	}
	that
	$\bar{I}$ is maximized at $a^* \approx 0.3164$
	with $\bar{I}(a^*) \approx 1.288$.
	%
	The constant $c_2$ is the square root of the latter.
\end{proof}

%

We can now state 
the
announced 
rate characterization for 
the Landweber iteration.

%


\begin{proposition}
\label{p:LW}
	Let $\{ x_k \}_{k \geq 0}$ be the Landweber iterates 
	generated by \eqref{e:LW}
	with 
	step size $\sigma > 0$ as in \eqref{e:tau}.
	%
	Let $\nu > 0$ and $r \geq 0$.
	%
	Set
	\begin{align}
	\label{e:LWE}
		\Delta_{\nu}^r(x^\dag - x_0)
		:=
		\sup_{k \geq 0}
		\varepsilon_{k}^{-(r + \nu)}
		\norm{ (T^* T)^r (x^\dag - x_k) }{}
		%
		\quad\text{where}\quad
		%
		\varepsilon_{k}
		:=
		\frac{1}{\sigma}
		\frac{r + \nu}{k + r + \nu}
		.
	\end{align}

	%
	Then, 
%	for $\normiii{\cdot}{\nu}$ defined in \eqref{e:normiii},
	%
	\begin{align}
	\label{e:p:LW}
% 		e^{-(p + r)}
% 		( \sigma \varepsilon(1) )^{r + \nu}
		c_1
		\sqrt{ \tfrac{\nu}{r + \nu} }
		\normiii{x^\dag - x_0}{\nu}
	\leq
		\Delta_{\nu}^r(x^\dag - x_0)
	\leq
		c_2
		\normiii{x^\dag - x_0}{\nu}
		,
	\end{align}
	where $c_2$ is from \eqref{e:c2} and
	\begin{align}
	\label{e:p:LW:c1}
		c_1
		:=
		\inf_{k \geq 0}
		\left\{
			(1 - \sigma \varepsilon_{k})^{k}
			\left[
				\varepsilon_{k+1} / \varepsilon_{k}
			\right]^{r + \nu}
		\right\}
		\geq
		( \sigma \varepsilon_1 / e)^{r + \nu}
		.
	\end{align}
\end{proposition}

%\begin{itemize}
%\item 
%	constants are independent of $x^\dag$
%\item
%	$\varepsilon_{k} \sim k^{-1}$ for $k \to \infty$.
%\item
%	if $0 < \nu < 1$,
%	meaningful choice $p = 1$ and $r = 1/2$.
%	Leads to well-known rate $\mathcal{O}(k^{-(\nu + 1/2)})$.
%\end{itemize}


\begin{proof}
	Fix $x^\dag \in X$.
	%
	We abbreviate
	$\mu := \mu_{x^\dag - x_0}$ 
	and
%	$\varepsilon := \varepsilon_{k}$,
%	as well as
	$\Delta := \Delta_{\nu}^r(x^\dag - x_0)$.
	%
	%
	The quantity $\varepsilon$ has been defined
	such as
	to satisfy
	\begin{align} 
		\varepsilon_{k}
		=
		\argmax_{\lambda \geq 0}
		f(\lambda)
		\quad\text{with}\quad
		f(\lambda) := \lambda^{2(r + \nu)} (1 - \sigma \lambda)^{2k}
		\quad
		\forall k \geq 0
		.
	\end{align}
	%
	%
	To enhance readability we will assume 
	for the remainder of the proof
	that $x_0 = 0$ and $\sigma = 1$.
	
	%
	
	For the first inequality of \eqref{e:p:LW},
	define $\mu^{2r}(A) := \int_A \lambda^{2 r} d \mu(\lambda)$
	on Borel subsets $A \subset [0, \infty)$.
	%
	Given $0 < \Lambda \leq \norm{T^* T}{}$,
	determine the integer $k \geq 0$ 
	by
	$\varepsilon_{k+1} < \Lambda \leq \varepsilon_{k}$.
	%
	Then
	%\begin{subequations}
	%\label{e:p:LW:mu2r1}
	\begin{align}
		\mu^{2 r}([0, \Lambda])
		\leq
		\mu^{2 r}([0, \varepsilon_{k}])
	& \leq
		(1 - \varepsilon_{k})^{-2 k}
		\int_{[0, \varepsilon_{k}]}
		\lambda^{2 r}
		(1 - \lambda)^{2 k}
		d\mu(\lambda)
	\\
	& \leq
		(1 - \varepsilon_{k})^{-2 k}
		\varepsilon_{k}^{2 (r + \nu)}
		\Delta^2
		\quad
		\text{(by definition of $\Delta$)}
	\\
	& \leq
		(1 - \varepsilon_{k})^{-2 k}
		\left[
			\varepsilon_{k} / \varepsilon_{k+1}
		\right]^{2 (r + \nu)}
		\Lambda^{2(r + \nu)}
		\Delta^2
	\\
	&
	\leq
		c_1^{-2} 
		\Lambda^{2(r + \nu)}
		\Delta^2
		.
	\end{align}
	%\end{subequations}
	%
	Since $\Lambda > 0$ was arbitrary,
	we obtain the bound
	\begin{align}
	\label{e:p:LW:mu2r}
		\normiii{\mu^{2 r}}{r + \nu}
	\leq 
		c_1^{-1}
		\Delta
		.
	\end{align}
	%
	The first inequality of \eqref{e:p:LW}
	now
	follows from:
	\begin{align}
		\normiii{\mu}{\nu}^2
%	=
%		\sup_{t > 0}
%		t^{-2 \nu}
%		\mu([0, t))
	& =
		\sup_{t > 0}
		t^{-2 \nu}
		\int_{[0, t)} 
		\lambda^{-2 r}
		d\mu^{2 r}(\lambda)
	\stackrel{\eqref{e:Tail1}}{\leq}
		\tfrac{r + \nu}{\nu}
		\normiii{\mu^{2 r}}{r + \nu}^2
	\stackrel{\eqref{e:p:LW:mu2r}}{\leq}
		\tfrac{r + \nu}{\nu}
		c_1^{-2}
		\Delta^2
		.
	\end{align}
	%
	
	%
	
	
	
	
	%
	
	We now prove the second inequality in \eqref{e:p:LW}.
	%
	In view of \eqref{e:tau},
	the measure $\mu$ is supported on $[0, 1/\sigma]$,
	and recall that we have assumed $\sigma = 1$.
	%
	So consider
	\begin{align}
		\varepsilon_{k}^{-2(r + \nu)}
		\norm{ (T^* T)^r (x^\dag - x_k) }{}^2
	& =
		\varepsilon_k^{-2(r + \nu)}
		\int_{[0, 1]} \lambda^{2 r} (1 - \lambda)^{2 k } d\mu(\lambda)
		.
	\intertext{%
		In the case $k = 0$,
		this is majorized by $\mu([0, 1]) \leq \normiii{\mu}{\nu}^2$.
		%
		Otherwise,
		defining $I(\lambda) := \mu([0, \lambda))$ 
		and integrating by parts as in 
		the proof of Lemma \ref{l:mu},
		we estimate
	}
	& =
		\varepsilon_k^{-2(r + \nu)}
		\int_{[0, 1]} (-(\lambda^{2 r} (1 - \lambda)^{2 k})' I(\lambda) d\lambda
	\\
	& \leq
		\varepsilon_k^{-2(r + \nu)}
		\int_{[0, 1]} 2 k \lambda^{2 r} (1 - \lambda)^{2 k - 1} I(\lambda) d\lambda
		.
	\end{align}
	%
	Estimating $I(\lambda) \leq \lambda^{2 \nu} \normiii{\mu}{\nu}^2$,
	and employing Lemma \ref{l:c2}
	with $s := 2k$ and $a := 2 (r + \nu)$,
	the second inequality in \eqref{e:p:LW} is obtained.
	%
	%
	This completes the proof of \eqref{e:p:LW}.
\end{proof}

%

Unfortunately, the gap between our ``constants'' in \eqref{e:p:LW}
is not robust in $\nu$.
%
For example, if $r = 0$,
we have
$c_2 / c_1 \sim e^{\nu}$ as $\nu \to \infty$.
%
%We believe this could be significantly improved.
%
Qualitatively though,
\eqref{e:p:LW} means that
it is necessary and sufficient 
for
the asymptotic convergence rate $\nu$
that the data be in $\IX_\nu$,
which explains the notoriously slow convergence
of the Landweber iteration.
%
In the limit $\nu \to \infty$,
the upper estimate in \eqref{e:p:LW}
yields the asymptotics
\begin{align}
	\norm{ (T^* T)^r (x^\dag - x_k) }{}
	\leq
	c_2
	e^{-k}
	(1 + \tfrac12 k^2 \nu^{-1} + \mathcal{O}(\nu^{-2}))
	\normiii{ x^\dag - x_0 }{\nu}
	\quad\text{as}\quad
	\nu \to \infty
	,
\end{align}
for any fixed iteration $k \geq 0$.
%
%
This bound indicates that, for a fixed iteration $k$,
only a limited amount of smoothness can be exploited
(cf.~\cite[Section 3, Example 3]{Mathe2004}).

%

From \cite[Theorem 1 and Corollary 1]{FlemmingHofmannMathe2011},
an estimate similar to \eqref{e:p:LW} with $r = 0$
can be obtained
if we assume
$d(r) \sim r^{-\nu/(1 - \nu)}$ for sufficiently large $r > 0$
for the distance function
(see Section \ref{s:i:dr}).

%

We believe it is natural for $\normiii{\cdot}{\nu}$
to appear in \eqref{e:p:LW}
rather than the interpolation norm $\norm{\cdot}{\nu : \gamma}$.
%
The following example illustrates this in 
the limiting case $\nu = \gamma$.
%
As in the $k$-fold Tikhonov regularization,
however,
it might be possible and meaningful
to relate the error of the $k$-th iterate 
to the $\norm{\cdot}{\nu : k}$ norm of the data.

\begin{example}
\label{x:IX:LW}
	Let $T$ and $x^\dag$ be as in Example \ref{x:IX}.
	%
	Recall that $\norm{x^\dag}{1} = \infty$,
	while $\normiii{x^\dag}{1} \approx 1.096$.
	%
	We will estimate
	$\Delta_{\nu}^{0}(x^\dag)$ for $\nu := 1$.
	%
	Since $\norm{T}{} = 1$, we set $\sigma := 1$ to satisfy \eqref{e:tau}.
	%
	Then, for any integer $k \geq 0$,
	\begin{align}
	\label{e:xLW}
		\varepsilon_k^{-2}
		\norm{ x^\dag - x_k }{}^2
		=
		\varepsilon_k^{-2}
		\int_{[0, \infty)}
		(1 - \lambda)^{2 k}
		d \mu_{x^\dag}
		=
		(k + 1)^2
		\sum_{n \geq 1}
		n^{-3}
		(1 - n^{-1})^{2 k}
		.
	\end{align}
	%
	Splitting the sum at $n = k$,
	and exploiting 
	in the finite subsum
	the fact that $t \mapsto t^{-3} (1 - t^{-1})^{2k}$
	peaks at $t = 1 + \tfrac23 k$
	we have
	\begin{align}
		\sum_{n \geq 1} (\cdots)
		\leq
		\sum_{n \leq k} (1 + \tfrac23 k)^{-3}
		+
		\sum_{n > k} n^{-3}
		\leq
		C
		(k + 1)^{-2}
		\quad
		\forall k \geq 1
	\end{align}
	for some $C \geq 0$ independent of $k$,
	meaning that \eqref{e:xLW}
	is bounded uniformly in $k \geq 0$.
	%
	We find numerically that 
	$|\Delta_{\nu}^{0}(x^\dag)|^2 = \eqref{e:xLW}|_{k = 0} \approx (1.096)^2$,
	$\eqref{e:xLW}|_{k = 1} \approx (0.5453)^2$,
	$\eqref{e:xLW}|_{k = 2} \approx (0.5475)^2$,
	and
	$\eqref{e:xLW}$
	is decreasing for $k \geq 2$
	with 
	$\lim_{k \to \infty} \eqref{e:xLW} = (1/2)^2$.
	%
	This limit can also be verified using \eqref{e:soo}.
	%
	The infimum in $c_1$ is achieved at $k = 2$,
	so $c_1 = 1/3$.
	%
	%
	Thus \eqref{e:p:LW} reads 
	$1/3 \times 1.096 \leq 1.096 \leq 1.135 \times 1.096$.
	%
	%
	% % Mathematica source:
	% % nu = 1; e[k_] := nu/(k + nu); 
	% % c1[k_] := Limit[(1 - e[kk])^kk (e[kk + 1]/e[kk])^nu, kk -> k]; 
	% % ListPlot[Table[{(e[1]/Exp[1])^nu, c1[k]}, {k, 0, 10}] // Transpose, PlotRange -> All];
	% % c1[2]
	%
	% % Values of sqrt{\eqref{e:xLW}}:
	% % Sqrt[ (0 + 1)^2 Sum[n^(-3) (1-1/n)^(2 0), {n, 1, infinity}] ]
	% % Sqrt[ (1 + 1)^2 Sum[n^(-3) (1-1/n)^(2 1), {n, 1, infinity}] ]
	% % Sqrt[ (2 + 1)^2 Sum[n^(-3) (1-1/n)^(2 2), {n, 1, infinity}] ]
	%
	%
	If we use any of the iterates $(n^{-3/2} (1 - n^{-1})^{k})_{n \geq 1}$ as 
	the initial guess $x_0$,
	then \eqref{e:p:LW} reads
	$1/3 \times 1/\sqrt{2} \leq 1/2 \leq 1.135 \times 1/\sqrt{2}$.
\end{example}

%

%

%
%A similar statement to \eqref{e:p:LW} may hold for $r > -\nu$
%but the proof does not go through without modification.
%
Our primary motivation for introducing $r \geq 0$ in \eqref{e:LWE}
was to estimate the number of iterations needed
for
the stopping rule based on the discrepancy principle \eqref{e:MDPk}.
%
We briefly comment on this.
%
Fix $\tau > 1$.
%
%
In view of \eqref{e:dade} and \eqref{e:tau}
we have
\begin{align}
	\left|
		\norm{y^\dag - T x_k}{}
		-
		\norm{y^\delta - T x_k^\delta }{}
	\right|
	\leq
	\norm{ (y^\dag - T x_k) - (y^\delta - T x_k^\delta) }{}
	\leq
	\delta
\end{align}
for each iteration $k \geq 0$.
%
Hence, if 
$\norm{ y^\dag - T x_k }{} \leq \delta (\tau - 1)$
then \eqref{e:MDPk} is satisfied.
%
Using \eqref{e:p:LW} with $r = 1/2$,
this is fulfilled
for any $k \geq 0$ such that
$ %\begin{align}
	\delta (\tau - 1)
	\leq
	{\varepsilon}_k^{\nu + 1/2}
	c_2
	\normiii{x^\dag - x_0}{\nu}
$. %\end{align}
%Here,
%$ %\begin{align}
%	{\varepsilon}_k
%	=
%	\frac{1}{\sigma}
%	\frac{\nu + 1/2}{k + \nu + 1/2}
%	.
%$ %\end{align}
%
%
This implies at most $k^\star \sim \delta^{-2 / (2 \nu + 1)}$ 
iterations
until \eqref{e:MDPk} is met,
as in \cite[Theorem 6.5]{EnglHankeNeubauer1996}
but with the slightly weaker assumption that 
the data 
% $x^\dag - x_0$
is in $\IX_\nu$ rather than in $X_\nu$.
%
%
To reproduce the conclusion of \cite[Theorem 6.5]{EnglHankeNeubauer1996},
we still need to verify that
the discrepancy principle \eqref{e:MDPk} implies
the error rate
$\norm{ x^\dag - x_k^\delta }{} \lesssim \delta^{2 \nu / (2 \nu + 1)}$.
%
Here we cannot use 
the inequality 
\cite[(4.66)]{EnglHankeNeubauer1996},
\begin{align}
	\norm{ x^\dag - x_k }{}
	\leq
	\norm{ x^\dag - x_0 }{\nu}^{1 / (2 \nu + 1)} 
	\norm{ y^\dag - T x_k }{}^{2 \nu / (2 \nu + 1)}
	,
\end{align}
as in \cite[Proof of Theorem 6.5]{EnglHankeNeubauer1996}
for
the corresponding error rate
with data in $X_\nu$.
%
Instead, we estimate
\begin{align}
	\norm{ x^\dag - x_k }{}^2
	& =
	\int_{[0, \infty)} (1 - \sigma \lambda)^{2k} d\mu(\lambda)
\\
	& 
	\leq
	\int_{[0, \varepsilon_{k^\star})} d\mu(\lambda)
	+
	\varepsilon_{k^\star}^{-1}
	\int_{[\varepsilon_{k^\star}, \infty)} \lambda (1 - \sigma \lambda)^{2k} d\mu(\lambda)
\\
	& 
	\leq
	\varepsilon_{k^\star}^{2 \nu}
	\normiii{\mu}{\nu}^2
	+
	\varepsilon_{k^\star}^{-1}
	\norm{ y^\dag - T x_k }{}^2
	%
	\lesssim
	\delta^{4 \nu / (2 \nu + 1)}
%\\
%	& \lesssim
%	\delta^{4 \nu / (2 \nu + 1)}
%	+
%	\delta^{-2 / (2 \nu + 1)}
%	\delta^2
%	=
%	2 \, \delta^{4 \nu / (2 \nu + 1)}
	,
\end{align}
%
and using this in the error splitting \eqref{e:LW:es} 
yields
$\norm{ x^\dag - x_k^\delta }{} \lesssim \delta^{2 \nu / (2 \nu + 1)}$.
%
%
% In particular, the inequality ``$\lesssim$''
% is valid in \eqref{e:qok}
% for the stopping rule given by the discrepancy principle \eqref{e:MDPk}.
% %
% We do not know whether ``$\gtrsim$'' also holds.

%



% %%%%%%%%%%%%%%%%%%%%%%%%%%%%%%%%%%%%%%%%%%%%%%%%%%%%%%%%%%%%%%%%%
% \section{Numerical experiments}
% \label{s:n}
% %%%%%%%%%%%%%%%%%%%%%%%%%%%%%%%%%%%%%%%%%%%%%%%%%%%%%%%%%%%%%%%%%


%%%%%%%%%%%%%%%%%%%%%%%%%%%%%%%%%%%%%%%%%%%%%%%%%%%%%%%%%%%%%%%%%
\section{Conclusions}
\label{s:99}
%%%%%%%%%%%%%%%%%%%%%%%%%%%%%%%%%%%%%%%%%%%%%%%%%%%%%%%%%%%%%%%%%

We have 
introduced and investigated
families of Banach spaces 
and
their interrelations:
the Hilbert scale $X_\nu$, the interpolation spaces $X_{\nu:\gamma}$,
and
the spaces $\IX_\nu = X_{\nu : 2 \nu}$.
%
We have shown that 
the interpolation spaces $X_{\nu:\gamma}$
are most adequate for the characterization of convergence rates 
in (iterated) Tikhonov regularization,
while
$\IX_\nu$ are better suited
for the Landweber iteration.
%
%
\REV{%
This insight should facilitate
the verification 
of 
the convergence rates.
%
For instance,
in many situations,
$X_{\nu:\gamma}$
can be understood as (a subspace of)
the
Besov space $B_{2,\infty}^s$, 
for which many characterizations are known.
}%\REV


%%%%%%%%%%%%%%%%%%%%%%%%%%%%%%%%%%%%%%%%%%%%%%%%%%%%%%%%%%%%%%%%%
\section*{Acknowledgments}
%%%%%%%%%%%%%%%%%%%%%%%%%%%%%%%%%%%%%%%%%%%%%%%%%%%%%%%%%%%%%%%%%

I thank 
J.~Flemming for providing me with a copy of \cite{Flemming2012};
M.~Hansen for the reference to the Gagliardo completion;
O.~Scherzer for general comments on the initial draft;
the anonymous referees for their comments,
in particular
concerning quasioptimality \cite{RausHamarik2007} and Besov spaces.
%
%
The note was mainly written while at RICAM, Austrian Academy of Sciences, Linz (AT).
%
%
Supported in part by ANR-12-MONU-0013.

%%%%%%%%%%%%%%%%%%%%%%%%%%%%%%%%%%%%%%%%%%%%%%%%%%%%%%%%%%%%%%%%%
\section*{References}
%%%%%%%%%%%%%%%%%%%%%%%%%%%%%%%%%%%%%%%%%%%%%%%%%%%%%%%%%%%%%%%%%

%%%%%%%%%%%%%%%%%%%%%%%%%%%%%%%%%%%%%%%%%%%%%%%%%%%%%%%%%%%%%%%%%%%%%%%%%%%%%%%%
\bibliographystyle{iopart-num}
\bibliography{refs}
%%%%%%%%%%%%%%%%%%%%%%%%%%%%%%%%%%%%%%%%%%%%%%%%%%%%%%%%%%%%%%%%%%%%%%%%%%%%%%%%


%%%%%%%%%%%%%%%%%%%%%%%%%%%%%%%%%%%%%%%%%%%%%%%%%%%%%%%%%%%%%%%%%%%%%%%%%%%%%%%%
\end{document}
%%%%%%%%%%%%%%%%%%%%%%%%%%%%%%%%%%%%%%%%%%%%%%%%%%%%%%%%%%%%%%%%%%%%%%%%%%%%%%%%
